\documentclass[12pt,a4paper]{report}


% ============================================================
% Packages
% ============================================================
\usepackage[utf8]{inputenc}
\usepackage[margin=1in]{geometry}
\usepackage{graphicx}
\usepackage{amsmath}
\usepackage{amssymb}
\usepackage{hyperref}
\usepackage{cite}
\usepackage{setspace}
\usepackage{fancyhdr}
\usepackage{titlesec}
\usepackage{enumitem}
\usepackage{booktabs}
\usepackage{longtable}

% ============================================================
% Page style
% ============================================================
\pagestyle{fancy}
\fancyhf{}
\cfoot{\thepage}
\renewcommand{\headrulewidth}{0pt}



% ============================================================
% Hyperref setup
% ============================================================
\hypersetup{
    colorlinks=true,
    linkcolor=blue,
    filecolor=magenta,
    urlcolor=cyan,
    citecolor=green,
}

% ============================================================
% Title / chapter formatting
% ============================================================
\titleformat{\chapter}[hang]
{\normalfont\huge\bfseries}{\thechapter.}{1em}{\Huge}
\titlespacing*{\chapter}{0pt}{-20pt}{20pt}

% ============================================================
% Document information (kept for \maketitle if needed)
% ============================================================
\title{AcadAI: An AI-Powered Research Paper Assistant\\
\large Minor Project Report}
\author{
    Aman Dhaubanjar (079BCT010)\\
    Anjal Satyal (079BCT014)\\
    Ankit Prasad Kisi (079BCT017)\\
    \\
    \textit{Artificial Intelligence and Machine Learning (AIML) Cluster}
}
\date{\today}

% ============================================================
\begin{document}
% ============================================================

% ============================================================
% TITLE PAGE
% ============================================================
\begin{titlepage}
    \centering
    \fontsize{14pt}{16pt}\selectfont
    \includegraphics[width=1.5in]{logo.png}\\[2.5mm]

    \textbf{\LARGE Tribhuvan University}\\
    \textbf{\Large Institute of Engineering}\\
    \textbf{\Large Pulchowk Campus}\\[1.7cm]

    {\Large \textbf{A MINOR PROJECT REPORT ON }}\\[0.2cm]
    {\Large AcadAI: An AI-Powered Research Paper Assistant}\\[2cm]

    {\bf Submitted By:}\\[0.2cm]
    \begin{tabular}{ll}
        Aman Dhaubanjar  & (079BCT010)\\
        Anjal Satyal     & (079BCT014)\\
        Ankit Prasad Kisi & (079BCT017)
    \end{tabular}\\[2cm]

    {\bf Submitted To: } \\[0.2cm]
    Department of Electronics and Computer Engineering \\
    IOE, Pulchowk Campus \\
    Lalitpur, Nepal \\ [2cm]

    {\medium 10 May, 2026}\par

    \vfill
\end{titlepage}

% ============================================================
% FRONT MATTER  –  Roman numerals
% ============================================================
\pagenumbering{roman}
\setcounter{page}{1}

% ------------------------------------------------------------
% PAGE OF APPROVAL
% ------------------------------------------------------------
\chapter*{Page of Approval}
\addcontentsline{toc}{chapter}{Page of Approval}

\begin{center}
    {\large \textbf{AcadAI: An AI-Powered Research Paper Assistant}}\\[0.5cm]
    {\large \textbf{Minor Project Report}}\\[1cm]
\end{center}

\noindent This project report has been prepared by the following students of the
Department of Electronics and Computer Engineering, Institute of Engineering,
Pulchowk Campus, Tribhuvan University, in partial fulfilment of the requirements
for the Bachelor's Degree in Computer Engineering.

\vspace{1cm}

\noindent \textbf{Submitted By:}\\[0.3cm]
\begin{tabular}{lll}
    Aman Dhaubanjar   & 079BCT010 & \underline{\hspace{5cm}} \\[0.6cm]
    Anjal Satyal      & 079BCT014 & \underline{\hspace{5cm}} \\[0.6cm]
    Ankit Prasad Kisi & 079BCT017 & \underline{\hspace{5cm}} \\
\end{tabular}

\vspace{1.5cm}

% TBD – Add supervisor name, department, and approval signatures when available
\noindent \textbf{Supervised By:}\\[0.3cm]
\begin{tabular}{ll}
    Name:       & \underline{\hspace{7cm}} \\[0.6cm]
    Designation: & \underline{\hspace{7cm}} \\[0.6cm]
    Signature:  & \underline{\hspace{7cm}} \\[0.6cm]
    Date:       & \underline{\hspace{7cm}} \\
\end{tabular}

\vspace{1.5cm}

\noindent \textbf{Department of Electronics and Computer Engineering}\\
Institute of Engineering, Pulchowk Campus\\
Tribhuvan University\\
Lalitpur, Nepal

\newpage

% ------------------------------------------------------------
% COPYRIGHT
% ------------------------------------------------------------
\chapter*{Copyright}
\addcontentsline{toc}{chapter}{Copyright}

\vspace*{3cm}
\begin{center}
    \textcopyright\ 2026 Aman Dhaubanjar, Anjal Satyal, Ankit Prasad Kisi
\end{center}

\vspace{1cm}
\noindent All rights reserved. This report is submitted in partial fulfilment of
the requirements for the Bachelor's Degree in Computer Engineering at the Institute
of Engineering, Pulchowk Campus, Tribhuvan University. No part of this report may
be reproduced, stored in a retrieval system, or transmitted in any form or by any
means—electronic, mechanical, photocopying, recording, or otherwise—without the
prior written permission of the authors.

\newpage

% ------------------------------------------------------------
% ACKNOWLEDGEMENTS
% ------------------------------------------------------------
\chapter*{Acknowledgements}
\addcontentsline{toc}{chapter}{Acknowledgements}

We would like to express our sincere gratitude to the faculty members of the
Department of Computer Engineering for their continuous support and guidance
throughout this project. Their expertise and insights have been invaluable in
shaping our understanding of advanced AI and NLP techniques.

We are also grateful to the open-source community for providing excellent tools
and libraries such as LangGraph, Docling, Qdrant, and Groq, which have been
instrumental in the development of this project.

Finally, we thank our families and peers for their constant encouragement and
understanding during the course of this work.

\vspace{1cm}
\begin{flushright}
\textit{Aman Dhaubanjar}\\
\textit{Anjal Satyal}\\
\textit{Ankit Prasad Kisi}
\end{flushright}

\newpage

% ------------------------------------------------------------
% ABSTRACT
% ------------------------------------------------------------
\chapter*{Abstract}
\addcontentsline{toc}{chapter}{Abstract}

Reading research papers can be overwhelming, especially for students and
early-stage researchers. The language is dense, the structure varies across
paper types, and identifying core contributions often takes more effort than it
should.

AcadAI is an AI-powered web application that makes reading research papers
simpler and more structured. It classifies each paper by type, generates a
personalised step-by-step reading guide, and answers comprehension questions
grounded in the specific section being read. At its core is a section-aware
hybrid retrieval system that combines semantic search and keyword matching,
constrained to the active reading section rather than searching the entire
document. Tables and figures are processed as distinct content types rather than
plain text, improving the quality of retrieved information.

The system is designed for undergraduate students, early-stage researchers, and
anyone exploring a new research area who wants structured guidance through
academic literature. Technologies used include FastAPI, React, LangGraph,
PostgreSQL, Qdrant, and Docling.

\newpage

% ------------------------------------------------------------
% TABLE OF CONTENTS
% ------------------------------------------------------------
\tableofcontents
\addcontentsline{toc}{chapter}{Contents}
\newpage

% ------------------------------------------------------------
% LIST OF FIGURES
% ------------------------------------------------------------
\listoffigures
\addcontentsline{toc}{chapter}{List of Figures}
\newpage

% ------------------------------------------------------------
% LIST OF TABLES
% ------------------------------------------------------------
\listoftables
\addcontentsline{toc}{chapter}{List of Tables}
\newpage

% ------------------------------------------------------------
% LIST OF ABBREVIATIONS
% ------------------------------------------------------------
\chapter*{List of Abbreviations}
\addcontentsline{toc}{chapter}{List of Abbreviations}

\begin{longtable}{ll}
\textbf{Abbreviation} & \textbf{Expansion} \\[0.3cm]
API   & Application Programming Interface \\
BM25  & Best Matching 25 \\
DB    & Database \\
LLM   & Large Language Model \\
MRR   & Mean Reciprocal Rank \\
NLP   & Natural Language Processing \\
OCR   & Optical Character Recognition \\
ORM   & Object-Relational Mapping \\
PDF   & Portable Document Format \\
QA    & Question Answering \\
RAG   & Retrieval-Augmented Generation \\
RRF   & Reciprocal Rank Fusion \\
SPA   & Single Page Application \\
TF-IDF & Term Frequency–Inverse Document Frequency \\
UI    & User Interface \\
\end{longtable}

\newpage

% ============================================================
% MAIN CONTENT  –  Arabic numerals
% ============================================================
\pagenumbering{arabic}
\setcounter{page}{1}

% ============================================================
% CHAPTER 1: INTRODUCTION
% ============================================================
\chapter{Introduction}

Reading research papers efficiently is a critical skill for students and
researchers, yet it remains challenging due to complex technical language, dense
content, and unstructured presentation of ideas. Beginners often struggle to
identify core contributions, navigate relevant sections, and follow the logical
flow of a paper, leading to inefficient reading patterns and poor comprehension.
This project addresses these issues by developing AcadAI, an AI-powered web
application that provides structured, step-by-step guidance for reading any
uploaded research paper.\\
AcadAI analyses the document, classifies it by paper type, extracts
content-type-aware chunks, and generates a personalised reading guide with
section-specific comprehension questions. At answer time, a section-scoped
hybrid retrieval pipeline retrieves only the most relevant chunks from the
section being read, grounding the generated answer in retrieved evidence rather
than the full document context. This separation of guide generation from answer
generation is a deliberate architectural decision that prevents the system from
reducing to a simple document injection wrapper.

% ------------------------------------------------------------
\section{Background}
% ------------------------------------------------------------

The exponential growth of scientific publications has created both opportunities
and challenges for the research community. While access to knowledge has never
been greater, the ability to efficiently navigate and comprehend academic
literature remains a critical bottleneck, particularly for students and
early-career researchers. Reading research papers requires not only domain
expertise but also the ability to identify key contributions, understand
methodological approaches, and synthesise complex ideas across multiple
sections.\\
Traditional approaches to reading research papers rely heavily on individual
experience and intuition. Students often struggle with identifying the most
relevant sections to read first, understanding technical terminology in context,
and maintaining a coherent mental model of the paper's contributions. This
cognitive overhead leads to inefficient reading patterns where readers may spend
excessive time on peripheral content while missing crucial insights.\\
Recent advances in Natural Language Processing and Large Language Models have
opened new possibilities for intelligent document understanding and assistance.
Retrieval-Augmented Generation (RAG) systems have demonstrated remarkable
capabilities in answering questions by retrieving relevant context from document
collections. However, existing tools typically perform global retrieval without
structural awareness, fail to adapt guidance to paper type, and treat tables and
figures as undifferentiated text rather than distinct semantic content types.
AcadAI addresses each of these gaps through targeted architectural decisions
described in the following chapters.

% ------------------------------------------------------------
\section{Problem Statement}
% ------------------------------------------------------------

Despite the growing importance of research literacy, many students and
early-stage researchers struggle to read and comprehend academic papers
effectively. Traditional reading methods and existing AI tools fail to guide
beginners through complex language, dense technical content, and varied paper
structures. Specific limitations include: the absence of paper-type-aware
reading guidance, global retrieval that returns chunks from irrelevant sections,
inability to process tables and figures as semantically meaningful units, and no
structured mechanism for section-specific comprehension assessment. These gaps
highlight the need for an intelligent system that can classify papers, structure
their content by type, and provide grounded, section-scoped question answering
to support comprehension.

% ------------------------------------------------------------
\section{Objectives}
% ------------------------------------------------------------

To design and implement a section-aware hybrid retrieval-augmented generation
system that classifies research papers by type, generates structured reading
guides, and answers section-specific comprehension questions using
content-type-aware chunking and metadata-filtered retrieval.

Specific objectives are:
\begin{itemize}
    \item Build a system to ingest research PDFs and extract text, tables, and
          figures while preserving section hierarchy.
    \item Classify each paper as Theory, Applied, or Survey using a trained
          classifier.
    \item Generate personalised, type-conditioned step-by-step reading guides.
    \item Support section-scoped hybrid retrieval and grounded question answering.
    \item Provide a web interface for browsing papers and interacting with the
          guided reading experience.
\end{itemize}

% ------------------------------------------------------------
\section{Scope}
% ------------------------------------------------------------

The scope of this project includes the design and implementation of a web-based
system for research paper analysis and guided reading. The system supports
section-aware document parsing using Docling, paper type classification using a
trained multinomial logistic regression model, content-type-aware chunking of
text, tables, and figures, hybrid RAG combining BM25 and dense embeddings with
Reciprocal Rank Fusion and cross-encoder reranking, section-scoped metadata
filtering in Qdrant, guide generation conditioned on paper type, and
section-specific question answering. The system is scoped to computer science
domain papers in PDF format.

\subsection*{Out of Scope}
The project will not address multi-paper comparative analysis or literature
review generation, automated citation network analysis, real-time collaborative
features, mobile application development, formula extraction, and non-English
language support.

% ------------------------------------------------------------
\section{Importance of the Project}
% ------------------------------------------------------------

AcadAI addresses a genuine gap in the tooling available to students and
researchers approaching unfamiliar academic literature. By combining paper type
classification, content-type-aware chunking, and section-scoped retrieval, the
system reduces the cognitive effort required to navigate a research paper and
provides structured, grounded guidance at the point of reading. For
undergraduate students undertaking a first literature review, or early-stage
researchers entering a new sub-field, such a tool can substantially reduce the
time required to identify key contributions and understand methodological
choices. The system also demonstrates a replicable architectural pattern for
building section-aware RAG applications over long structured documents.

% ------------------------------------------------------------
\section{Chapter Summary}
% ------------------------------------------------------------

The remainder of this report is organised as follows. Chapter~2 surveys related
work and the theoretical background underpinning the system. Chapter~3 details
the requirements and methodology. Chapter~4 presents the system design and
architecture. Chapter~5 describes the implementation module by module. Chapter~6
covers testing and evaluation. Chapter~7 presents and discusses the results.
Chapter~8 concludes the report and outlines future work.

% ============================================================
% CHAPTER 2: LITERATURE REVIEW
% ============================================================
\chapter{Literature Review}

% ------------------------------------------------------------
\section{Research Paper Analysis Tools}
% ------------------------------------------------------------

The landscape of tools designed to assist readers of academic literature has
evolved considerably in recent years, shifting from simple keyword-based search
engines toward AI-mediated comprehension systems that actively interpret,
summarise, and contextualise scholarly content. The tools surveyed below
represent the current state of the art and motivate the specific design
decisions made in \textit{AcadAI}.

\subsubsection{Semantic Scholar}

Semantic Scholar~\cite{semanticscholar} is a freely accessible, AI-driven
academic search engine developed by the Allen Institute for AI. Indexing over
200~million papers across disciplines, the platform employs machine learning to
surface influential citations, map connections between research works, and
generate concise automated summaries. Its \textit{TL;DR} feature distils the
key contribution of a paper into a single sentence, enabling rapid triage of
large result sets. Despite these strengths, Semantic Scholar is primarily a
discovery and search platform: it assists researchers in \emph{finding} relevant
papers rather than in \emph{deeply understanding} the content of a single
document. The summaries it provides are global rather than section-specific, and
the tool does not offer a structured, interactive reading roadmap for an
individual paper.

\begin{figure}[h]
    \centering
    \includegraphics[width=1\textwidth]{schemanticscholar.png}
    \caption{Semantic Scholar}
    \label{fig:semanticscholar}
\end{figure}

\subsubsection{Explainpaper}

Explainpaper~\cite{explainpaper} is an AI-powered reading assistant that
simplifies dense academic papers by allowing users to upload a document,
highlight passages of interest, and receive plain-English explanations generated
on demand. While Explainpaper effectively lowers the cognitive barrier to
individual passages, it operates on flat, unstructured text and does not exploit
the hierarchical organisation inherent in academic papers. It does not provide
structural guidance on how to navigate a paper, nor does it prioritise sections
according to a reader's specific information needs, and it performs no
section-aware retrieval.

\begin{figure}[h]
    \centering
    \includegraphics[width=1\textwidth]{explainpaper.png}
    \caption{Explainpaper}
    \label{fig:explainpaper}
\end{figure}
% ------------------------------------------------------------
\section{PDF Extraction and OCR}
% ------------------------------------------------------------

Research documents are predominantly distributed in the Portable 
Document Format (PDF), which presents significant challenges for 
automated text extraction. Unlike plain text formats, PDF encodes 
content as a stream of graphical operators rather than semantic 
text units, meaning the logical reading order, paragraph boundaries, 
and structural elements must be reconstructed algorithmically 
\cite{pdfminer2023}.

\subsection{Text-Based PDF Extraction}

For digitally authored PDFs, text extraction tools operate directly 
on the embedded character streams. Tools such as PDFMiner 
\cite{pdfminer2023}, PyMuPDF \cite{pymupdf2023}, and pdfplumber 
\cite{pdfplumber2023} retrieve character-level glyphs along with 
positional metadata including bounding box coordinates, font size, 
and font weight. These positional attributes are critical for 
downstream tasks such as heading detection and column 
disambiguation in multi-column academic layouts.

However, raw character extraction is insufficient for academic 
document understanding. Research papers typically contain complex 
layouts with two-column formats, mathematical equations, figures, 
tables, and footnotes interspersed with body text. Nave 
concatenation of extracted characters produces malformed text 
that corrupts downstream NLP tasks \cite{vila2021}. This problem 
motivated the development of layout-aware extraction frameworks.

Document layout analysis approaches the problem as a computer 
vision task, treating page images as inputs and identifying 
semantic regions such as titles, paragraphs, figures, and tables 
using object detection models.Layoutparser \cite{shen2021} 
introduced a unified framework combining layout detection with 
downstream text extraction. Building on this line of work, 
Docling \cite{docling2024} extends layout-aware extraction with 
a full document understanding pipeline. Docling applies a 
TableFormer model \cite{nassar2022} for structured table 
reconstruction and a dedicated layout model for region 
classification, producing a rich hierarchical document 
representation that preserves reading order, table cell 
structure, figure boundaries, and provenance metadata including 
per-element page numbers and bounding box coordinates.

The provenance data exposed by Docling, specifically the 
\texttt{prov} attribute on each document item containing 
\texttt{page\_no} and \texttt{bbox} fields, enables fine-grained 
location tracking of every extracted text element. This 
capability is particularly valuable for citation-level navigation 
in reading assistance systems, where retrieved text chunks must 
be mappable back to specific page regions in the original document.

% After explaining layout-aware extraction, before OCR subsection
\begin{figure}[h]
    \centering
    % \includegraphics[width=0.85\textwidth]{figures/docling_pipeline.png}
    \caption{Docling document understanding pipeline showing layout 
    detection, table structure recognition, and provenance extraction 
    stages for a two-column academic PDF.}
    \label{fig:docling_pipeline}
\end{figure}

\subsection{Scanned PDFs and OCR}

A significant proportion of academic documents, particularly 
older publications and conference proceedings, exist as scanned 
image PDFs with no embedded text layer. Optical Character 
Recognition (OCR) is required to recover textual content from 
these documents. Traditional OCR engines such as Tesseract 
\cite{smith2007} apply image binarisation, line segmentation, 
and character classification pipelines to produce text output, 
though accuracy degrades substantially on low-resolution scans, 
two-column layouts, and mathematical notation.

Deep learning approaches have substantially advanced OCR 
accuracy. PaddleOCR \cite{du2020} applies a detection-then-recognition 
pipeline using differentiable binarisation for text region 
detection and a scene text recognition model for character 
transcription, achieving competitive performance on document 
benchmarks. EasyOCR \cite{easyocr2021} similarly applies a 
two-stage deep learning architecture and supports over eighty 
languages with a straightforward Python interface.

Docling integrates multiple OCR backends including Tesseract, 
EasyOCR, and RapidOCR as interchangeable engines, applying them 
selectively on pages with low text density detected during the 
initial conversion pass. This hybrid strategy, processing 
digitally authored pages with the native text layer while 
falling back to OCR for image-heavy or scanned pages, optimises 
both accuracy and throughput for heterogeneous document 
collections.

\subsection{GPU Acceleration}

Modern document understanding pipelines apply neural models for 
layout detection, table structure recognition, and OCR, each 
requiring significant computation per page. GPU acceleration 
is therefore essential for acceptable throughput in interactive 
systems. Docling exposes a configurable accelerator device 
parameter that routes tensor operations to CUDA-enabled GPUs 
when available \cite{docling2024}. For a typical thirty-page 
academic paper, GPU-accelerated conversion completes in 
approximately seventeen to nineteen seconds compared to 
several minutes on CPU, making real-time upload processing 
feasible.

% ------------------------------------------------------------
\section{Metadata Extraction}
% ------------------------------------------------------------

Structured metadata extraction from academic papers encompasses 
the identification of bibliographic fields including title, 
authors, abstract, keywords, section headings, and citation 
references. Accurate metadata extraction underpins document 
indexing, retrieval, categorisation, and reading guide 
generation.

\subsection{Heuristic and Rule-Based Approaches}

Early metadata extraction systems applied hand-crafted rules 
exploiting typographic conventions in academic publishing. 
Titles typically appear as the largest font on the first page, 
abstracts follow a labelled heading, and author affiliations 
appear in smaller font between the title and abstract. 
Systems such as ParsCit \cite{councill2008} applied 
conditional random fields trained on positional and 
typographic features to extract citations and header fields 
from research papers.

GROBID \cite{lopez2009} advanced this approach by applying 
a cascade of sequence labelling models to segment and label 
document zones at multiple granularities, from document-level 
header fields down to individual citation components. GROBID 
remains a widely used baseline for academic metadata extraction, 
particularly for citation parsing, though its accuracy on 
non-standard layouts and documents from non-Western publishers 
is limited \cite{tkaczyk2018}.

\subsection{LLM-Assisted Extraction}

Large language models have emerged as a powerful alternative 
to rule-based metadata extraction, particularly for fields 
requiring semantic understanding rather than positional 
heuristics. Prompting an LLM with the first page text of a 
document and requesting structured field extraction leverages 
the model's understanding of academic writing conventions to 
identify title, abstract, and author fields even in 
non-standard layouts \cite{wei2022}.

The approach is particularly effective when combined with 
structured output constraints. Using JSON schema-constrained 
generation, LLMs can be prompted to return metadata as a 
typed dictionary with explicit field names, enabling 
programmatic consumption of extracted values without fragile 
string parsing \cite{openai2023}. Confidence scoring can be 
derived from field coverage, where the proportion of 
successfully extracted non-null fields indicates extraction 
quality.

A practical hybrid strategy applies layout-based extraction 
as the primary path, using heading detection to identify the 
abstract region and section names from the document structure, 
and reserves LLM-based extraction as a fallback for documents 
where the primary path achieves insufficient field coverage. 
This combination achieves near-complete metadata coverage 
across diverse document types while minimising LLM API calls 
for well-structured documents \cite{docling2024}.

\subsection{Section Heading Extraction}

Section heading identification is a prerequisite for both 
metadata extraction and document chunking. Headings are 
characterised by elevated font size, bold weight, or 
numbered prefix patterns such as Arabic numerals, Roman 
numerals, or alphabetic labels. Vila et al. \cite{vila2021} 
demonstrated that token-level classification combining 
textual and positional features substantially outperforms 
purely positional heuristics for heading detection in 
scientific documents.

Docling exposes extracted headings with their hierarchical 
level inferred from font size and numbering depth, providing 
a structured list of section titles with associated page 
numbers that forms the input to section hierarchy detection.

% ------------------------------------------------------------
\section{Section Detection and Hierarchy Extraction}
% ------------------------------------------------------------

The logical structure of an academic paper is organised as a 
hierarchy of sections, subsections, and nested content. 
Recovering this hierarchy from extracted text is essential 
for reading comprehension support, targeted retrieval, and 
coherent chunking, as section boundaries define the natural 
semantic units of a research document.

\subsection{Why Section Hierarchy Matters}

Reading comprehension research has long established that 
hierarchical document structure supports comprehension by 
providing readers with a mental map of content organisation 
\cite{kintsch1978}. A three-pass reading strategy, in which 
readers first survey section titles and abstracts, then read 
methodology and results sections in depth, and finally 
synthesise conclusions, is a widely recommended approach for 
efficiently extracting understanding from dense research papers 
\cite{keshav2007}. Automated reading guides that follow this 
strategy must therefore ground their instructions in the 
actual section hierarchy of the target document.

For information retrieval, section-aware chunking produces 
chunks that respect semantic boundaries rather than splitting 
content arbitrarily at fixed token counts. Retrieval systems 
that filter by section allow users to direct questions at 
specific document regions, improving precision when the user 
already has structural knowledge of the paper. Formal 
evaluation of section-aware retrieval has demonstrated 
improved answer relevance compared to section-agnostic 
chunking on scientific QA benchmarks \cite{yang2023}.

\subsection{Section Segmentation Approaches}

Section segmentation can be approached as a sequence 
labelling task, a classification problem at the text block 
level, or a rule-based post-processing step applied to 
layout-extracted headings.

Rule-based approaches exploit numbering conventions and 
font-weight signals. A section heading numbered \texttt{3.2} 
is unambiguously a level-two subsection under section three. 
Regular expression matching on common academic numbering 
patterns combined with font size thresholding provides a 
reliable baseline for well-formatted papers \cite{mao2012}.

Machine learning approaches treat section boundary detection 
as a binary classification task at each text block, using 
features including relative font size, vertical position, 
line spacing, and preceding content. DocBank \cite{li2020} 
provides a large-scale benchmark of annotated academic 
document structures that has been used to train layout 
classification models. SciPDF Parser and similar tools 
apply trained models to assign semantic labels including 
title, abstract, section heading, paragraph, figure caption, 
and table to each extracted text block.

\subsection{Hierarchy Tree Construction}

Once section headings are identified with their associated 
heading levels, a section hierarchy tree is constructed by 
interpreting level assignments as parent-child relationships. 
A level-one heading opens a new top-level section; a 
subsequent level-two heading becomes a child of the most 
recently opened level-one section. This stack-based 
construction algorithm is robust to typical academic 
numbering schemes.

Practical challenges in hierarchy construction include 
inconsistent heading level assignment by layout models, 
unnumbered sections such as Abstract and Acknowledgements 
that must be inferred from position and content, and 
appendices that follow the main body with separate numbering. 
Post-processing heuristics normalise these cases by 
assigning unnumbered headings to the appropriate level based 
on positional context and known heading vocabulary.

The output of section hierarchy extraction is a structured 
tree where each node carries the section title, inferred 
heading level, page number of first occurrence, and 
references to child sections. This tree is serialised as a 
JSON artifact that serves as the primary input to the 
document chunking pipeline, ensuring that chunk boundaries 
align with semantic section boundaries rather than arbitrary 
token counts.

\subsection{Confidence and Coverage Metrics}

Evaluating the quality of extracted section hierarchies 
requires comparing detected sections against ground truth 
annotations. In the absence of ground truth for a given 
document, proxy metrics such as heading count, maximum 
detected depth, and coverage of known section vocabulary 
(Introduction, Related Work, Methodology, Experiments, 
Conclusion) provide a confidence estimate. A high confidence 
score indicates that the extracted hierarchy is likely 
complete and correctly structured, while a low score 
triggers fallback handling such as coarser chunking or 
guide generation with reduced section specificity.
% ------------------------------------------------------------
\section{Retrieval-Augmented Generation and Question Answering}
% ------------------------------------------------------------

\subsection{Retrieval-Augmented Generation}

Retrieval-Augmented Generation (RAG)~\cite{lewis2020} combines information
retrieval with neural text generation to enable large language models to produce
responses grounded in external knowledge sources. Instead of relying solely on
the model's parametric memory, RAG dynamically retrieves relevant documents from
a knowledge base at inference time, allowing the system to incorporate
domain-specific information.

\begin{figure}[h]
    \centering
    \includegraphics[width=0.8\textwidth]{rag.png}
    \caption{RAG Pipeline}
    \label{fig:rag}
\end{figure}

RAG systems operate in two main stages:

\begin{enumerate}
    \item \textbf{Retrieval:} Given a query $q$, the retriever selects the
    top-$k$ most relevant documents or passages from a corpus:
    \begin{equation}
    D_{\text{retrieved}} = \text{TopK}(\{d_1, d_2, \ldots, d_N\}, q, k)
    \end{equation}
    The retriever can be based on sparse methods (e.g., BM25), dense vector
    similarity, or a hybrid combination of both.

    \item \textbf{Generation:} The language model conditions its response on
    both the query and the retrieved context to generate a grounded answer:
    \begin{equation}
    P(a|q) = \text{LLM}(q, D_{\text{retrieved}})
    \end{equation}
    This grounding reduces hallucinations and improves factual consistency by
    constraining the model to the retrieved evidence.
\end{enumerate}

By decoupling knowledge storage from the language model parameters, RAG systems
enable scalable knowledge updates, improved interpretability through retrieved
evidence, and better performance on knowledge-intensive tasks such as question
answering and document understanding.

\subsection{Information Retrieval Fundamentals}

\subsubsection{BM25 Scoring Function}

BM25 (Best Matching 25) is a ranking function used by search engines to estimate
the relevance of documents to a query. Given a query $Q$ containing keywords
$q_1, \ldots, q_n$, the BM25 score of a document $D$ is:

\begin{equation}
\text{score}(D, Q) = \sum_{i=1}^{n} \text{IDF}(q_i) \cdot
\frac{f(q_i, D) \cdot (k_1 + 1)}{f(q_i, D) + k_1 \cdot
\left(1 - b + b \cdot \frac{|D|}{\text{avgdl}}\right)}
\end{equation}

where $f(q_i, D)$ is the frequency of term $q_i$ in document $D$, $|D|$ is the
length of document $D$, $\text{avgdl}$ is the average document length in the
collection, $k_1$ and $b$ are tuning parameters (typically $k_1 = 1.2$ and
$b = 0.75$), and $\text{IDF}(q_i)$ is the inverse document frequency of term
$q_i$.

The inverse document frequency is calculated as:

\begin{equation}
\text{IDF}(q_i) = \ln\left(\frac{N - n(q_i) + 0.5}{n(q_i) + 0.5} + 1\right)
\end{equation}

where $N$ is the total number of documents and $n(q_i)$ is the number of
documents containing $q_i$.

\subsubsection{Dense Vector Retrieval}

Dense retrieval uses neural networks to encode queries and documents into
continuous vector spaces. A typical dense retrieval system consists of:

\begin{enumerate}
    \item \textbf{Encoder Network:} A neural network (often BERT-based) that
    maps text to fixed-dimensional vectors:
    \begin{equation}
    \mathbf{v}_q = \text{Encoder}_Q(q), \quad
    \mathbf{v}_d = \text{Encoder}_D(d)
    \end{equation}

    \item \textbf{Similarity Function:} Typically cosine similarity or dot
    product:
    \begin{equation}
    \text{sim}(q, d) = \frac{\mathbf{v}_q \cdot \mathbf{v}_d}
    {|\mathbf{v}_q| \cdot |\mathbf{v}_d|}
    \end{equation}

    \item \textbf{Training Objective:} Contrastive learning to maximise
    similarity between relevant query-document pairs and minimise similarity
    with negative examples.
\end{enumerate}

\subsubsection{Hybrid Retrieval and Reciprocal Rank Fusion}

Naive hybrid retrieval combines sparse and dense scores through weighted
aggregation:

\begin{equation}
\text{score}_{\text{hybrid}}(d, q) = \alpha \cdot \text{score}_{\text{BM25}}
(d, q) + (1-\alpha) \cdot \text{score}_{\text{dense}}(d, q)
\end{equation}

where $\alpha \in [0, 1]$ controls the balance between sparse and dense
retrieval. However, this approach is sensitive to score scale differences
between BM25 and dense retrieval, which occupy incomparable numerical ranges.
AcadAI instead uses Reciprocal Rank Fusion (RRF)~\cite{cormack2009}, a
rank-based fusion method that is robust to these scale differences. For a
document $d$ appearing at rank $r$ in a given ranked list, its RRF contribution
is:

\begin{equation}
\text{RRF}(d, r) = \frac{1}{k + r}
\end{equation}

where $k=60$ is a smoothing constant. The final RRF score is the sum of
contributions across all ranked lists, making it robust to outlier scores in
either retrieval system.

\begin{figure}[h]
    \centering
    \includegraphics[width=0.8\textwidth]{hybrid_rag.png}
    \caption{Hybrid Retrieval}
    \label{fig:hybridrag}
\end{figure}

\subsubsection{Cross-Encoder Reranking}

Cross-encoder models differ from bi-encoders in that they jointly encode the
query and a candidate document in a single forward pass, allowing the model to
attend to interactions between query tokens and document tokens. This produces
substantially more accurate relevance scores than bi-encoder similarity but at
the cost of $O(n)$ forward passes per query. In practice, cross-encoders are
applied as second-stage rerankers over a small candidate set produced by the
first-stage retriever. AcadAI uses \texttt{ms-marco-MiniLM-L-12-v2}, a
cross-encoder fine-tuned on the MS MARCO passage ranking dataset, as its
reranker.

\subsection{Section-Aware Metadata Filtering}

Section-aware retrieval incorporates document structure by filtering chunks
based on metadata. Given a query $q$ and the active reading section $s$,
retrieval in AcadAI applies a hard filter rather than a soft weight:

\begin{equation}
D_{\text{filtered}} = \{d \in D \mid \text{section\_id}(d) = s\}
\end{equation}

\begin{equation}
\text{score}_{\text{section}}(d, q, s) = \text{score}(d, q) \cdot
\mathbf{1}[\text{section\_id}(d) = s]
\end{equation}

This hard filter is implemented as a Qdrant payload filter applied before
approximate nearest neighbour search, constraining the retrieval corpus to
chunks belonging to the active section. This reduces the effective search space
by an order of magnitude compared to global retrieval and is the primary
mechanism by which AcadAI achieves section-specific grounding.

\subsection{Multinomial Logistic Regression for Text Classification}

Multinomial logistic regression is a generalisation of binary logistic
regression to multi-class classification problems. Given an input feature vector
$\mathbf{x}$ and $K$ classes, the model computes the probability of class $k$
as:

\begin{equation}
P(y = k \mid \mathbf{x}) = \frac{e^{\mathbf{w}_k^T \mathbf{x} + b_k}}
{\sum_{j=1}^{K} e^{\mathbf{w}_j^T \mathbf{x} + b_j}}
\end{equation}

where $\mathbf{w}_k$ and $b_k$ are the weight vector and bias for class $k$. In
AcadAI, $\mathbf{x}$ is a TF-IDF feature vector computed from the concatenated
title and abstract, and $K = 3$ corresponding to Theory, Applied, and Survey
paper types. The model is trained on a labelled dataset of arXiv paper metadata
and survey paper repository data described in Chapter~\ref{ch:reqmeth}.

% ------------------------------------------------------------
\section{Workflow Orchestration and LangGraph}
% ------------------------------------------------------------

% TBD – Content to be added: explain workflow graphs and state-based
% orchestration; describe multi-step pipelines for ingestion, classification,
% indexing, and QA using LangGraph.

% ------------------------------------------------------------
\section{Related Research Works}
% ------------------------------------------------------------

Several prior works explore techniques relevant to this project.
\textit{Adaptive Search Query Generation and Refinement in Systematic Literature
Review}~\cite{badami2023} investigates adaptive query generation strategies to
improve document retrieval, emphasising iterative query refinement.
\textit{Corrective Retrieval-Augmented Generation}~\cite{yan2024} introduces
corrective retrieval mechanisms that evaluate and refine retrieved evidence
before answer generation. \textit{FAIR-RAG}~\cite{asl2024fair} proposes a
feedback-driven retrieval process where the system re-queries documents when
confidence is low. \textit{HySemRAG}~\cite{godinez2025} presents a hybrid
retrieval framework combining semantic search and keyword filtering for
automated literature synthesis.

Although these works do not aim to generate structured reading guidelines for
research papers, their techniques in adaptive query refinement, corrective
retrieval, and hybrid retrieval are directly relevant. AcadAI builds upon these
ideas while adding paper type classification and content-type-aware chunking as
novel contributions specific to the research paper reading domain.

% ------------------------------------------------------------
\section{Gap Analysis}
% ------------------------------------------------------------

While the tools discussed above make important contributions, each addresses
only a subset of the challenges faced by a student attempting to read an
unfamiliar paper. Explainpaper provides localised explanations but ignores
document structure. Semantic Scholar excels at discovery but does not guide deep
reading. FAIR-RAG demonstrates the value of iterative retrieval but is not
designed with the anatomy of research papers in mind, and none of these tools
adapt their behaviour to whether a paper is theoretical, applied, or a survey.

\textit{AcadAI} bridges these gaps through a section-aware hybrid RAG
architecture with paper type classification. Every chunk is tagged with
section-level metadata enabling retrieval to be constrained to the active
reading section. A multinomial classifier trained on arXiv and survey repository
data ensures that guide structure adapts to paper type. Content-type-aware
chunking ensures that tables and figures contribute meaningful semantic
representations rather than noisy markdown artifacts.

% ------------------------------------------------------------
\section{Chapter Summary}
% ------------------------------------------------------------

This chapter surveyed the existing tools and theoretical foundations that
motivate the design of AcadAI. The gap analysis confirms that no existing tool
combines paper-type-aware guide generation, content-type-aware chunking, and
section-scoped hybrid retrieval. The theoretical background in RAG, BM25, dense
retrieval, RRF, and cross-encoder reranking provides the conceptual basis for
the system described in subsequent chapters.

% ============================================================
% CHAPTER 3: REQUIREMENTS
% ============================================================
\chapter{Requirements}
\label{ch:reqmeth}

% ------------------------------------------------------------
\section{Feasibility Study}
% ------------------------------------------------------------

\subsection{Technical Feasibility}
The system is built on established technologies and frameworks proven effective
in similar applications. Docling provides robust PDF parsing with layout
preservation and section detection. The BAAI/bge-small-en-v1.5 model provides
384-dimensional dense embeddings suitable for semantic similarity search. Qdrant
supports both dense vector representations and sparse BM25 representations
within a unified hybrid retrieval framework. Groq provides fast inference for
\texttt{llama-3.3-70b-versatile} for guide and answer generation,
\texttt{llama-3.1-8b-instant} for table summarisation, and
\texttt{llama-4-scout-17b-16e-instruct} for multimodal figure summarisation.
LangGraph enables complex multi-step AI workflows as structured state machines.
All components run on consumer-grade hardware without GPU acceleration.

\subsection{Scheduling Feasibility}
The project is divided into well-defined phases. Core components including
document ingestion, paper type classification, hybrid retrieval, guide
generation, and question answering have been implemented. Remaining work
includes stabilising retrieval quality through chunk size tuning, completing the
PostgreSQL integration for rich media storage, building the three-panel web
interface, and conducting the evaluation described in Chapter~\ref{ch:testing}.

\subsection{Resource Feasibility}
The project requires a laptop or desktop with 8--16~GB RAM. All required
software libraries are freely available and open-source. API access uses the
Groq free tier (30 requests per minute) and the Qdrant free cloud tier. The
team consists of three members with complementary skills in backend development,
NLP, and frontend development.

% ------------------------------------------------------------
\section{Functional Requirements}
% ------------------------------------------------------------

\begin{itemize}
    \item The system accepts PDF files of research papers and extracts text,
          tables, and figures while preserving section hierarchy using Docling.
    \item The system classifies each paper as Theory, Applied, or Survey using
          a trained multinomial logistic regression model operating on title and
          abstract.
    \item The system generates text, table, and figure chunks as distinct content
          types. Table chunks are summarised using an LLM. Figure chunks are
          summarised using a multimodal LLM with the image file as input.
    \item The system indexes all chunks into Qdrant with dense and sparse
          vectors, with each chunk carrying a \texttt{section\_id},
          \texttt{section\_title}, \texttt{content\_type}, and
          \texttt{image\_path} (for figure chunks) in its payload.
    \item The system stores rich media including figure file paths, table
          summaries, and technical terms in PostgreSQL keyed by
          \texttt{section\_id}.
    \item The system generates a structured reading guide conditioned on paper
          type with each guide step linked to a specific section and carrying
          \texttt{needs\_figures} and \texttt{needs\_tables} flags.
    \item The system generates section-specific comprehension questions for each
          guide step.
    \item The system retrieves relevant chunks for a user question using
          section-scoped hybrid retrieval, applies RRF fusion, cross-encoder
          reranking, relevance threshold filtering, and near-duplicate
          deduplication before generating an answer.
    \item The web interface presents the reading guide on the left panel, the
          paper PDF in the centre panel, and contextual components including
          answers, figures, tables, and technical terms in the right panel.
\end{itemize}

% ------------------------------------------------------------
\section{Non-Functional Requirements}
% ------------------------------------------------------------

Document parsing completes within 2 minutes for typical papers. Query response
time does not exceed 10 seconds under normal operating conditions. The system
handles papers up to 50 pages. Answers are grounded exclusively in retrieved
chunks, with no reliance on the language model's parametric knowledge for
paper-specific claims.

% ------------------------------------------------------------
\section{Tools and Technologies}
% ------------------------------------------------------------

\begin{table}[h]
\centering
\begin{tabular}{ll}
\toprule
\textbf{Component} & \textbf{Library / Service} \\
\midrule
Document parsing        & Docling \\
Pipeline orchestration  & LangGraph \\
Dense embedding         & BAAI/bge-small-en-v1.5 \\
Sparse retrieval        & BM25SparseEncoder \\
Vector store            & Qdrant \\
Cross-encoder reranker  & ms-marco-MiniLM-L-12-v2 \\
Paper classifier        & Multinomial logistic regression (scikit-learn) \\
Guide / QA LLM          & llama-3.3-70b-versatile (Groq) \\
Table summarisation     & llama-3.1-8b-instant (Groq) \\
Figure summarisation    & llama-4-scout-17b-16e-instruct (Groq) \\
Relational database     & PostgreSQL \\
Backend framework       & FastAPI \\
Frontend framework      & React \\
RAG evaluation          & RAGAS \\
\bottomrule
\end{tabular}
\caption{Software and libraries used in AcadAI}
\label{tab:software}
\end{table}

% ============================================================
\chapter{System Design}
% ============================================================

This chapter presents the complete system design of AcadAI. It covers the
high-level system architecture, the identified actors and use cases, the
relational database schema, the security design, and the interaction and
behavioural models of the system expressed through sequence, activity, and
workflow diagrams. The design decisions described here form the structural
foundation on which the implementation described in subsequent chapters is
built.

% ------------------------------------------------------------
\section{System Overview}
% ------------------------------------------------------------

AcadAI is an AI-assisted academic paper reading assistant that accepts
academic PDF documents as input, automatically extracts structured content,
classifies the paper by type, and generates a personalised multi-pass reading
guide. During active reading, the system supports section-scoped
question-answering and technical-term definition through a retrieval-augmented
generation (RAG) pipeline.

The system is composed of five principal subsystems that communicate in a
layered fashion: a React-based single-page application (SPA) serving the
user interface; a FastAPI backend providing the HTTP API; a LangGraph
orchestration layer managing AI workflows; a PostgreSQL relational database
for structured storage; and a Qdrant vector database for semantic retrieval.
External dependencies include the Groq API for large language model (LLM)
inference, HuggingFace-hosted sentence-transformer models for dense encoding,
and Docling for layout-aware PDF parsing.

\begin{figure}[h]
    \centering
    \includegraphics[width=\textwidth]{system_architecture_block.png}
    \caption{High-Level System Architecture}
    \label{fig:architecture}
\end{figure}

The top-level data flow is as follows. The user interacts with the React
frontend, which communicates with the FastAPI backend over HTTP. On PDF
upload, the backend triggers an asynchronous extraction pipeline that
processes the document through Docling, builds the section hierarchy, and
persists structured artefacts to PostgreSQL. A LangGraph workflow then
classifies the paper and generates a reading guide using the Groq-hosted Llama
model. Concurrently, the document is chunked and indexed into Qdrant for
retrieval. During a reading session, user questions are handled by a hybrid
retrieval pipeline that queries Qdrant, reranks results, and passes the
retrieved context to the LLM to generate a grounded answer.

% ------------------------------------------------------------
\section{Architectural Layers}
% ------------------------------------------------------------

The system follows a five-layer architecture with clear separation of
responsibilities across each layer. Each layer depends only on the layers
below it, which keeps the design modular and independently testable.

\subsection{Presentation Layer}

The presentation layer is a React~18 single-page application built with
TypeScript and Vite. It renders a three-pane workspace comprising a left
navigation panel (\texttt{PaperNavigation.tsx}), a central PDF viewer
(\texttt{PaperViewer.tsx}), and a right AI tools panel
(\texttt{AIToolsPanel.tsx}). The application uses TanStack Query for
server-state caching and background refresh, React Router for client-side
routing, react-pdf for in-browser PDF rendering, and Tailwind CSS with Radix
UI primitives for layout and interaction. Authentication tokens and the
selected paper identifier are persisted in browser \texttt{localStorage} to
maintain state across page refreshes.

\subsection{API Layer}

The API layer is a FastAPI service implemented in \texttt{backend/api/app.py}.
It exposes sixteen HTTP endpoints grouped into four functional categories:
authentication (\texttt{/api/auth/}), paper management
(\texttt{/api/papers/}), reading and AI interaction
(\texttt{/api/papers/\{id\}/}), and content management
(\texttt{/api/cms/papers/}). The service uses Pydantic models for request and
response validation and Uvicorn as the ASGI server. Upload endpoints schedule
extraction as a background task and expose a server-sent events (SSE) progress
stream so the frontend can display live processing status.

\subsection{Orchestration Layer}

The orchestration layer is implemented as a LangGraph state machine defined in
\texttt{backend/rag/graph.py}. The workflow graph contains five nodes:
\texttt{categorizer\_node}, \texttt{applied\_guide\_node},
\texttt{theoretical\_guide\_node}, \texttt{survey\_guide\_node}, and
\texttt{retrieve\_and\_qa\_node}. A conditional edge,
\texttt{route\_after\_categorizer}, branches execution based on the predicted
paper type. The shared state object, \texttt{AgentState}, carries the paper
title, abstract, full text, document identifier, paper type, generated guide,
and question-section pairs throughout the workflow.

\subsection{Persistence Layer}

The persistence layer is built on PostgreSQL with SQLAlchemy~2.0 as the ORM
framework. The schema, defined in
\texttt{backend/extraction/persistence/postgres\_store.py}, comprises thirteen
tables covering users, papers, sections, text blocks, tables, images,
references, reading guides, generated questions, and four junction tables
linking extracted elements to their parent sections. The
\texttt{PostgresPaperStore} class encapsulates all database operations,
including schema initialisation, extraction persistence with deduplication,
user management, and access control.

\subsection{Retrieval Layer}

The retrieval layer, implemented in
\texttt{backend/rag/retrieval/pipeline.py}, provides hybrid vector and sparse
retrieval over indexed document chunks. It combines the
\texttt{BAAI/bge-small-en-v1.5} sentence-transformer model (384-dimensional
dense embeddings) with a per-document BM25 sparse index. Qdrant serves as the
vector store for cosine similarity search, and FlashRank provides
cross-encoder reranking using the \texttt{ms-marco-MiniLM-L-12-v2} model.
The pipeline exposes two primary methods: \texttt{index()}, which chunks,
encodes, and uploads document content to Qdrant; and \texttt{query()}, which
retrieves, fuses, and reranks results for a given question.

% ------------------------------------------------------------
\section{Use Case Design}
% ------------------------------------------------------------

\subsection{Actors}

The system involves four actors:

\begin{itemize}
    \item \textbf{Registered User} — An authenticated user who can upload
    papers, manage their reading sessions, ask questions, and request
    AI-generated content.

    \item \textbf{System (Background Processor)} — An internal actor
    representing the asynchronous pipeline that runs extraction, guide
    generation, and Qdrant indexing after a user uploads a PDF.

    \item \textbf{External LLM Service (Groq)} — Receives generation prompts
    from the backend and returns structured text responses for guide
    generation, question answering, and term definition.

    \item \textbf{Vector Store (Qdrant)} — Receives index upserts and query
    vectors from the retrieval pipeline and returns ranked candidate chunks.
\end{itemize}

\subsection{Use Case Diagram}

\begin{figure}[h]
    \centering
    \includegraphics[width=\textwidth]{use_case_diagram.png}
    \caption{Use Case Diagram — AcadAI}
    \label{fig:usecase}
\end{figure}

\subsection{Use Case Descriptions}

The principal use cases for the Registered User are summarised in
Table~\ref{tab:usecases}.

\begin{longtable}{|p{1.2cm}|p{3.5cm}|p{9.3cm}|}
\hline
\textbf{UC ID} & \textbf{Use Case} & \textbf{Description} \\
\hline
\endfirsthead
\hline
\textbf{UC ID} & \textbf{Use Case} & \textbf{Description} \\
\hline
\endhead
UC01 & User Registration & User creates an account with email, password, and
optional display name. The password is hashed with PBKDF2-HMAC-SHA256 before
storage. \\
\hline
UC02 & User Login & User authenticates with email and password. On success,
the backend returns an HMAC-signed bearer token valid for 24 hours. \\
\hline
UC03 & Retrieve Current User & User calls \texttt{/api/auth/me} with a bearer
token to retrieve their profile (id, email, display name, created timestamp). \\
\hline
UC04 & Upload PDF Paper & User uploads a PDF file (maximum 50~MB). The backend
validates the file, creates a pending paper record, links it to the user, and
triggers the asynchronous extraction pipeline. \\
\hline
UC05 & Monitor Upload Progress & User subscribes to the SSE stream at
\texttt{/api/papers/\{id\}/progress} to receive phase-by-phase status updates
during extraction, guide generation, and Qdrant indexing. \\
\hline
UC06 & List Papers & User retrieves all papers linked to their account via
\texttt{/api/papers} (paginated, default limit 100). \\
\hline
UC07 & Retrieve Paper Bundle & User calls \texttt{/api/papers/\{id\}/bundle}
to fetch the full extracted paper bundle, including sections, text blocks,
tables, images, and references. \\
\hline
UC08 & Download PDF File & User retrieves the original uploaded PDF via
\texttt{/api/papers/\{id\}/pdf}. \\
\hline
UC09 & Retrieve Reading Guide & User calls \texttt{/api/papers/\{id\}/guide}
to fetch the category-specific three-pass reading guide. \\
\hline
UC10 & List Reading Questions & User retrieves pre-generated questions for a
paper via \texttt{/api/papers/\{id\}/questions}, which returns each question
with its generation status. \\
\hline
UC11 & Ask Question During Reading & User posts a question to
\texttt{/api/papers/\{id\}/chat} with an optional section filter. The backend
performs hybrid retrieval and returns a grounded LLM-generated answer with
source references. \\
\hline
UC12 & Generate Answer for Question & User requests answer generation for a
specific pre-generated question via
\texttt{/api/papers/\{id\}/questions/\{qid\}/generate}. The backend invokes
the retrieve-and-QA node and stores the result. \\
\hline
UC13 & Request Technical Term Definition & User calls
\texttt{/api/papers/\{id\}/technical-terms/generate} to receive AI-generated
definitions for technical terms identified in the paper. \\
\hline
UC14 & Delete Paper & User removes a paper via the CMS endpoint, which deletes
the paper record, associated extracted elements, and indexed vectors from
Qdrant. \\
\hline
UC15 & Navigate Section Hierarchy & User browses the collapsible section tree
in the left navigation panel to move through the paper structure. \\
\hline
UC16 & Synchronise PDF View with Section & User selects a section in the
navigation panel; the central PDF viewer scrolls to and highlights the
corresponding page range. \\
\hline
\caption{Use Case Descriptions — Registered User}
\label{tab:usecases}
\end{longtable}

The system-level use cases, performed automatically after a user uploads a
paper, are as follows. The extraction pipeline validates the PDF, parses its
content with Docling, builds the section hierarchy, and persists all extracted
elements to PostgreSQL (UC17--UC23). The LangGraph workflow then classifies
the paper type and generates the reading guide and associated questions
(UC24--UC27). Finally, the retrieval pipeline chunks the document, encodes
the chunks, and indexes them into Qdrant (UC26).

% ------------------------------------------------------------
\section{Database Design}
% ------------------------------------------------------------

\subsection{Entity-Relationship Model}

The relational schema is implemented in PostgreSQL and managed through
SQLAlchemy~2.0. It contains thirteen tables organised around the central
\texttt{papers} entity. Figure~\ref{fig:erd} presents the entity-relationship
diagram.

\begin{figure}[h]
    \centering
    \includegraphics[width=\textwidth]{erd_diagram.png}
    \caption{Entity-Relationship Diagram — PostgreSQL Schema}
    \label{fig:erd}
\end{figure}

\subsection{Table Descriptions}

\paragraph{papers.}
The core entity storing each uploaded academic paper. Key columns include
\texttt{paper\_name} (unique, case-insensitive), \texttt{title},
\texttt{abstract}, \texttt{pdf\_hash} (used for deduplication),
\texttt{source\_pdf\_path}, \texttt{document\_uuid} (links to file-system
artefacts), and \texttt{metadata\_json} (JSONB field storing authors, dates,
and page count). A partial unique index on \texttt{pdf\_hash} and a
functional index on \texttt{LOWER(paper\_name)} enforce deduplication at the
database level.

\paragraph{users.}
Stores user accounts with \texttt{email} (case-insensitive unique),
\texttt{display\_name}, \texttt{password\_hash}, and \texttt{password\_salt}.
The hash is produced by PBKDF2-HMAC-SHA256 with 200{,}000 iterations; the
salt is a 16-byte random value stored as a hexadecimal string.

\paragraph{user\_papers.}
A many-to-many junction table linking users to the papers they own. The
\texttt{user\_id} and \texttt{paper\_id} foreign keys both cascade on delete.
All protected API routes query this table to enforce per-user access control
before returning paper data.

\paragraph{sections.}
Stores the hierarchical section decomposition of each paper. Each row carries
a \texttt{section\_key} (e.g., \texttt{1.2.3}), \texttt{parent\_section\_id}
(self-referential FK, nullable for root sections), \texttt{level},
\texttt{page\_start}, \texttt{position}, and a \texttt{stats\_json} field
containing word and element counts for that section.

\paragraph{text\_blocks.}
Stores extracted paragraph-level text with \texttt{element\_id} (from
Docling), \texttt{page\_number}, \texttt{text\_content}, and
\texttt{metadata\_json} (bounding box, font, language).

\paragraph{tables\_data.}
Stores extracted table content in two forms: \texttt{markdown\_content} (for
display and retrieval) and \texttt{text\_content} (plain text). Both are
indexed and searchable.

\paragraph{images.}
Stores extracted figures with \texttt{image\_path} (pointing to a saved PNG
file), \texttt{caption}, and \texttt{metadata\_json} (bounding box,
dimensions, colour space).

\paragraph{references\_data.}
Stores parsed bibliographic references with \texttt{reference\_text} and
\texttt{metadata\_json} containing structured fields (authors, title, venue,
year, DOI).

\paragraph{paper\_guides.}
Stores the generated reading guide for each paper. Contains
\texttt{guide\_json} (the three-pass guide structure),
\texttt{guide\_plan\_json} (generation metadata), and
\texttt{question\_section\_pairs\_json} (a flattened list of question-section
mappings).

\paragraph{paper\_questions.}
Stores individual reading questions generated during the LangGraph workflow.
Each row holds \texttt{question\_text}, \texttt{section\_title},
\texttt{retrieval\_payload\_json} (pre-retrieved context chunks),
\texttt{answer\_text} (populated on demand), and a \texttt{status} field
(\texttt{pending}, \texttt{running}, \texttt{completed}, or
\texttt{failed}).

\paragraph{Junction tables.}
Three many-to-many junction tables — \texttt{section\_text\_blocks},
\texttt{section\_tables}, and \texttt{section\_images} — link extracted
elements to the sections they belong to, enabling the frontend to reconstruct
a section-scoped paper bundle.

\subsection{Key Relationships}

Table~\ref{tab:relationships} summarises the principal foreign key
relationships in the schema.

\begin{table}[h]
\centering
\begin{tabular}{|l|l|l|}
\hline
\textbf{Table} & \textbf{References} & \textbf{Relationship} \\
\hline
user\_papers & users, papers & Many-to-many (access control) \\
\hline
sections & papers, sections (self) & One paper to many sections; \\
 & & self-referential parent-child \\
\hline
text\_blocks & papers & One paper to many text blocks \\
\hline
tables\_data & papers & One paper to many tables \\
\hline
images & papers & One paper to many images \\
\hline
references\_data & papers & One paper to many references \\
\hline
paper\_guides & papers & One paper to one guide \\
\hline
paper\_questions & papers & One paper to many questions \\
\hline
section\_text\_blocks & sections, text\_blocks & Junction (many-to-many) \\
\hline
section\_tables & sections, tables\_data & Junction (many-to-many) \\
\hline
section\_images & sections, images & Junction (many-to-many) \\
\hline
\end{tabular}
\caption{Principal Foreign Key Relationships}
\label{tab:relationships}
\end{table}

% ------------------------------------------------------------
\section{Security Design}
% ------------------------------------------------------------

\subsection{Password Storage}

Passwords are never stored in plaintext. During registration, the backend
generates a 16-byte cryptographically random salt using
\texttt{secrets.token\_bytes(16)} and derives the stored hash using
PBKDF2-HMAC-SHA256 with 200{,}000 iterations. Both the hexadecimal hash and
the hexadecimal salt are stored in the \texttt{users} table. During login,
the submitted password is re-derived with the stored salt and compared against
the stored hash using \texttt{hmac.compare\_digest()}, a constant-time
comparison that prevents timing-based attacks.

\subsection{Session Management and Token Design}

After successful authentication, the backend issues a stateless bearer token
constructed as follows. A JSON payload containing the user identifier
(\texttt{uid}), email (\texttt{em}), and expiration timestamp (\texttt{exp},
set to 86{,}400 seconds from issue time) is base64url-encoded to produce the
message component. An HMAC-SHA256 signature is computed over this message
using the \texttt{AUTH\_SECRET\_KEY} environment variable and
base64url-encoded to produce the signature component. The token is formatted
as \texttt{<message>.<signature>}. On each protected request, the backend
decodes the message, recomputes the signature, compares it with the received
signature using \texttt{hmac.compare\_digest()}, and verifies that the
expiration timestamp has not elapsed.

\subsection{Access Control}

All paper-related API routes enforce ownership verification. Before returning
any paper data, the backend calls
\texttt{store.user\_has\_access\_to\_paper(user\_id, paper\_id)}, which
queries the \texttt{user\_papers} junction table. If no matching record
exists, the endpoint returns \texttt{403 Forbidden}. This server-side
enforcement ensures that a valid token for one user cannot be used to access
another user's papers.

\subsection{Token Transmission}

The frontend stores the token in browser \texttt{localStorage} under the key
\texttt{auth\_token} and attaches it to every authenticated request via the
\texttt{Authorization: Bearer <token>} HTTP header. The backend reads this
header on protected routes, decodes the token, and resolves the authenticated
user before processing the request.

% ------------------------------------------------------------
\section{Sequence Diagrams}
% ------------------------------------------------------------

\subsection{User Registration and Login}

Figure~\ref{fig:seq_auth} shows the interaction between the user, the React
frontend, the FastAPI backend, and PostgreSQL during account registration and
login.

\begin{figure}[h]
    \centering
    \includegraphics[width=\textwidth]{sequence_auth.png}
    \caption{Sequence Diagram — User Registration and Login}
    \label{fig:seq_auth}
\end{figure}

During registration, the user submits an email, password, and optional display
name. The frontend posts this payload to \texttt{/api/auth/register}. The
backend validates the inputs, generates a password salt and PBKDF2 hash,
inserts a \texttt{UserRecord} into PostgreSQL, creates an HMAC-signed bearer
token, and returns the token with the user profile. The frontend stores the
token in \texttt{localStorage} and navigates to the main workspace.

During login, the user submits credentials to \texttt{/api/auth/login}. The
backend retrieves the user record by email, re-derives the hash from the
submitted password and stored salt, compares it using a constant-time digest,
and, on a successful match, issues a new bearer token in the same format as
registration.

\subsection{PDF Upload and Offline Indexing}

Figure~\ref{fig:seq_upload} shows the interaction between the user, the
frontend, the FastAPI backend, the extraction pipeline, the LangGraph
workflow, PostgreSQL, Qdrant, and the Groq API during PDF upload and
background processing.

\begin{figure}[h]
    \centering
    \includegraphics[width=\textwidth]{sequence_upload.png}
    \caption{Sequence Diagram — PDF Upload and Offline Indexing}
    \label{fig:seq_upload}
\end{figure}

The user selects a PDF file in the frontend, which posts it as multipart form
data to \texttt{/api/papers/upload}. The backend validates the file, computes
its hash, creates a pending \texttt{PaperRecord}, links the authenticated user
to it via \texttt{user\_papers}, and schedules the extraction pipeline as a
background task. The endpoint returns immediately with the paper identifier
and a \texttt{pending} status so the frontend can begin polling the SSE
progress stream.

In the background, the extraction pipeline runs Docling on the PDF to produce
text blocks, tables, images, and references. Metadata extraction and section
hierarchy construction follow, and all results are persisted to PostgreSQL. A
progress event is emitted after each phase. The LangGraph workflow then
classifies the paper, generates the reading guide via the Groq API, and stores
the guide in \texttt{paper\_guides}. Once the guide is ready (approximately
25 seconds after upload), the frontend receives a progress event and fetches
the guide. Finally, the retrieval pipeline chunks the document, encodes the
chunks, and upserts them to Qdrant. When indexing completes (approximately
99 seconds after upload), the QA chat becomes active.

\subsection{Online Query — User Asks a Question}

Figure~\ref{fig:seq_qa} shows the interaction during a reading session when
the user submits a question in the chat panel.

\begin{figure}[h]
    \centering
    \includegraphics[width=\textwidth]{sequence_qa.png}
    \caption{Sequence Diagram — Online Query Phase}
    \label{fig:seq_qa}
\end{figure}

The user types a question in \texttt{ChatAssistant.tsx}, optionally selecting
a section filter from the dropdown. The frontend posts the message history and
optional section filter to \texttt{/api/papers/\{id\}/chat}. The backend
invokes the retrieval pipeline, which encodes the query with
\texttt{BAAI/bge-small-en-v1.5}, performs dense search in Qdrant, performs
BM25 sparse search, fuses the ranked lists using Reciprocal Rank Fusion
(RRF), reranks the top candidates with FlashRank, and returns the top chunks.
The backend constructs a prompt from the retrieved chunks and the conversation
history (last ten turns) and sends it to the Groq API. The LLM response, along
with source metadata (section title and page start), is returned to the
frontend, which renders the answer as formatted Markdown and displays clickable
source references. Clicking a source box scrolls the PDF viewer to the
relevant page.

\subsection{Reading Guide Retrieval}

After a paper is fully processed, the frontend calls
\texttt{/api/papers/\{id\}/bundle} and \texttt{/api/papers/\{id\}/guide} in
parallel. The bundle endpoint returns all extracted elements, which populate
the section tree and the InsightExtractor panel. The guide endpoint returns
the \texttt{guide\_json} from \texttt{paper\_guides}, which the
\texttt{PaperNavigation} component renders as three tabbed passes. Each pass
step lists the sections to read and the associated questions. Clicking a
question sends it to the chat panel as a pre-filled input.


% ------------------------------------------------------------
\section{LangGraph Workflow Design}
% ------------------------------------------------------------

Figure~\ref{fig:langgraph} shows the LangGraph state machine that orchestrates
the guide generation and question-answering workflows.

\begin{figure}[h]
    \centering
    \includegraphics[width=\textwidth]{langgraph_workflow.png}
    \caption{LangGraph Workflow — Guide Generation and QA Graphs}
    \label{fig:langgraph}
\end{figure}

The guide generation graph begins at the \texttt{categorizer\_node}, which
invokes the \texttt{TfidfPaperCategorizer} on the paper title and abstract to
produce a paper type label (\texttt{applied}, \texttt{theoretical}, or
\texttt{survey}). The conditional edge \texttt{route\_after\_categorizer}
routes execution to one of three guide nodes. Each guide node invokes
\texttt{llama-3.3-70b-versatile} via the Groq API with a category-specific
prompt. Applied paper prompts focus on problem statement, methodology, and
evaluation; theoretical paper prompts focus on theorems, proof strategies, and
formal structure; survey paper prompts focus on taxonomy, coverage gaps, and
research landscape. All three routes converge at a post-processing step that
normalises the guide structure, validates coverage, and ensures that each pass
step contains at least one question. The final guide is returned as a
structured \texttt{guide\_json} object.

The QA chat graph is invoked on each user question. Its single compound node,
\texttt{retrieve\_and\_qa\_node}, executes the full retrieval pipeline — dense
retrieval, BM25 sparse retrieval, RRF fusion, FlashRank reranking, threshold
filtering, and deduplication — and then calls the LLM with the assembled
context to produce a structured Markdown answer with source citations.

% ------------------------------------------------------------
\section{Module Design}
% ------------------------------------------------------------

Figure~\ref{fig:modules} shows the internal module structure of the frontend
and backend subsystems and their interactions.

\begin{figure}[h]
    \centering
    \includegraphics[width=\textwidth,height=0.9\textheight,keepaspectratio]{module_design.png}
    \caption{Module Design — Frontend and Backend Components}
    \label{fig:modules}
\end{figure}

On the frontend, \texttt{Index.tsx} acts as the application shell and
orchestrates state across the three panels. \texttt{PaperNavigation.tsx}
manages the reading guide tabs, section tree, upload form, and delete
controls. \texttt{PaperViewer.tsx} wraps the react-pdf renderer and handles
\texttt{scrollToSection()} and \texttt{jumpToPage()} commands triggered by
section selection or chat source clicks. \texttt{ChatAssistant.tsx} manages
the message history, section filter state, and API communication for the chat
endpoint. \texttt{AIToolsPanel.tsx} serves as the container for the chat and
insights panels. \texttt{InsightExtractor.tsx} displays extracted tables,
figures, and technical term definitions.

On the backend, the extraction pipeline (\texttt{extraction.py}) coordinates
the three extraction stages — ingestion, metadata, and section hierarchy —
and returns a structured \texttt{ExtractionOutput}. The persistence store
(\texttt{postgres\_store.py}) handles all database reads and writes behind a
clean class interface. The retrieval pipeline
(\texttt{rag/retrieval/pipeline.py}) encapsulates the full hybrid retrieval
logic. The LangGraph graph (\texttt{rag/graph.py}) wires the five workflow
nodes together, and \texttt{tfidf\_categorizer.py} loads and applies the
pre-trained classification model.

% ------------------------------------------------------------
\section{Data Flow Design}
% ------------------------------------------------------------

Figure~\ref{fig:dataflow} presents the end-to-end data flow through the
system, from PDF upload to an active question-answering session.

\begin{figure}[h]
    \centering
    \includegraphics[width=\textwidth,height=0.9\textheight,keepaspectratio]{data_flow.png}
    \caption{Data Flow Diagram — End-to-End Pipeline}
    \label{fig:dataflow}
\end{figure}

A PDF uploaded by the user enters the Docling conversion stage, which uses
its \texttt{StandardPdfPipeline} to parse text items with provenance
information (page number and bounding box). The resulting
\texttt{DoclingDocument} object is passed to the metadata extractor, which
produces a \texttt{ProcessedDocument} containing the title, abstract,
section list, and figure and table counts. The section hierarchy detector then
assigns hierarchical keys and parent-child relationships, writing the result
to \texttt{<document\_uuid>\_hierarchy.json} on disk. PostgreSQL persistence
inserts the paper, section, text block, table, image, and reference records
from this output.

In parallel with persistence, the TF-IDF classifier reads the title and
abstract and routes to the appropriate LangGraph guide node. The Groq LLM
produces a structured three-pass guide, which is stored in
\texttt{paper\_guides} and becomes visible in the frontend approximately 25
seconds after upload. Simultaneously, the section chunker splits the hierarchy
JSON into fine-grain (150-character) and coarse-grain (400-character) chunks.
The BM25 encoder is fitted on the document corpus and the BGE dense encoder
produces 384-dimensional embeddings for each chunk. Both are upserted to
Qdrant as hybrid points with section and page metadata in the payload. The QA
chat becomes active approximately 99 seconds after upload.

When the user asks a question, the query is encoded by the BGE model and sent
to Qdrant for dense retrieval, while BM25 scores are computed locally. RRF
fusion merges the two ranked lists, FlashRank reranks the top candidates, and
the filtered top-four chunks are passed to the Groq LLM as context. The LLM
returns a structured Markdown answer that the frontend renders alongside
clickable section-and-page source references.

% ------------------------------------------------------------
\section{Chapter Summary}
% ------------------------------------------------------------

This chapter presented the complete system design of AcadAI. The layered
architecture separates the React frontend, the FastAPI API layer, the LangGraph
orchestration layer, the PostgreSQL persistence layer, and the Qdrant retrieval
layer, each with well-defined responsibilities and interfaces. The use case
model identifies sixteen user-facing interactions and eleven system-level
automated processes. The relational database schema comprises thirteen tables
with cascade-delete referential integrity and JSONB columns for
semi-structured extraction output. The security design employs
PBKDF2-HMAC-SHA256 password hashing, HMAC-signed stateless bearer tokens, and
server-side ownership verification on every protected route. The sequence,
activity, and workflow diagrams presented in this chapter provide the
behavioural specification from which the implementation described in the
following chapter is derived.

% ============================================================
\chapter{Methodology}
% ============================================================

This chapter describes the methodology adopted in the design, development, and
evaluation of AcadAI. It covers the development approach, system architecture,
technology stack, backend and frontend implementation strategies, database
design, authentication mechanism, API layer design, testing approach, and
deployment configuration.

% ------------------------------------------------------------
\section{Development Methodology}
% ------------------------------------------------------------

The system was developed using an iterative incremental methodology. Rather
than attempting to implement the entire application in a single pass, the
project was structured as a sequence of tightly scoped vertical slices: PDF
ingestion, metadata extraction, section hierarchy construction, database
persistence, retrieval and reading-guide generation, web presentation, and user
interaction features including authentication, chat, and technical-term support.
Each increment introduced a testable, independently operable workflow that
integrated cleanly into the overall pipeline.

This methodology is reflected in the repository structure itself, which
separates extraction, retrieval, persistence, API delivery, and frontend
rendering into dedicated modules. Each stage could therefore be developed,
validated, and refined without destabilising adjacent components. The project
also employs dedicated unit tests, smoke tests, and integration tests to verify
correct behaviour at each increment before the next stage is introduced.

% ------------------------------------------------------------
\section{System Architecture}
% ------------------------------------------------------------

The implemented architecture follows a layered monorepo model with clear
separation of concerns across five functional layers. The frontend is a
single-page React application that communicates with a FastAPI backend over
HTTP. The backend orchestrates PDF processing, language-model inference,
retrieval, and persistence. PostgreSQL serves as the relational database for
structured storage, while Qdrant provides vector-based retrieval for
paper-level question answering and reading-guide generation.

The five architectural layers are as follows:

\begin{itemize}
    \item \textbf{Presentation layer}: A React-based user interface that
    supports authentication, paper browsing, PDF rendering, reading-guide
    navigation, and AI-assisted chat interaction.

    \item \textbf{API layer}: FastAPI endpoints that expose authentication,
    paper management, file upload, retrieval, question generation, and PDF
    delivery services.

    \item \textbf{Orchestration layer}: A LangGraph state machine that manages
    the sequencing of document processing stages and reasoning tasks, including
    classification, guide generation, and retrieval-augmented question
    answering.

    \item \textbf{Persistence layer}: A PostgreSQL-backed relational schema
    storing papers, sections, text blocks, tables, figures, user accounts,
    reading guides, and generated questions.

    \item \textbf{Retrieval layer}: A Qdrant-backed hybrid retrieval subsystem
    combining dense semantic embeddings, BM25 sparse retrieval, Reciprocal
    Rank Fusion, and cross-encoder reranking.
\end{itemize}

This architecture was selected because it supports both structured relational
storage and flexible document-centric metadata while keeping AI-assisted
processing modular and independently replaceable.

% ------------------------------------------------------------
\section{Technology Stack}
% ------------------------------------------------------------

The project employs a polyglot stack centred on Python for backend intelligence
and TypeScript with React for the user interface.

\subsection{Backend Technologies}

The backend is implemented with FastAPI and Uvicorn, providing a lightweight
asynchronous HTTP service with typed request and response models via Pydantic.
SQLAlchemy and psycopg are used for PostgreSQL access. Document ingestion and
structure extraction rely on Docling, while Groq-hosted large language models
are used for metadata fallback, guide generation, summarisation, question
answering, and technical-term definition. LangGraph provides deterministic
workflow orchestration, and the LangChain integration layer is used to invoke
Groq models. The retrieval stack combines Qdrant, a BM25 sparse encoder,
sentence-transformers for dense embeddings, and FlashRank for cross-encoder
reranking.

\subsection{Frontend Technologies}

The frontend is built with React~18 and Vite, providing fast development
iteration and optimised build output. React Router handles client-side route
management, TanStack Query manages server-state caching and background refresh,
and react-pdf renders PDF documents directly in the browser. The interface is
composed using Tailwind CSS and Radix-based UI primitives. TypeScript is used
throughout to enforce type safety across API contracts and component props.

\subsection{Data and Infrastructure Technologies}

PostgreSQL is the primary database, chosen for its support of relational
linking, transactional consistency, and JSONB columns for semi-structured
extraction results. Qdrant serves as the vector store because the project
requires similarity search over chunked document content with payload filtering
by section scope. Local file storage is used for extraction artefacts, generated
guides, and cached retrieval assets. Runtime configuration is managed through
environment variables with optional \texttt{.env} file loading.

% ------------------------------------------------------------
\section{Two-Phase Pipeline Design}
% ------------------------------------------------------------

AcadAI is organised as a two-phase processing pipeline orchestrated by a
LangGraph state machine.

\begin{figure}[h]
    \centering
    \begin{minipage}{0.48\textwidth}
        \centering
        \includegraphics[width=\textwidth]{offline_indexing_phase.jpg}
        \caption{Offline Indexing Phase}
        \label{fig:offline}
    \end{minipage}
    \hfill
    \begin{minipage}{0.48\textwidth}
        \centering
        \includegraphics[width=\textwidth]{online_query_phase.jpg}
        \caption{Online Query Phase}
        \label{fig:online}
    \end{minipage}
\end{figure}

\subsection{Offline Indexing Phase}

The offline indexing phase is triggered when a user uploads a PDF. The pipeline
traverses the following stages in sequence. Docling parsing produces a
hierarchy JSON, a full-text artefact, and a complete document JSON. The paper
type classifier reads the title and abstract and assigns one of three
categories: Theory, Applied, or Survey. The content-type-aware chunker then
processes text, table, and figure elements separately, applying large language
model (LLM) summarisation to table and figure chunks. All chunks are indexed
into Qdrant using dual dense and sparse vectors with section-level metadata,
while rich media references are written to PostgreSQL. Finally, reading-guide
generation runs conditioned on the classification result, and question
generation runs per guide step.

\subsection{Online Query Phase}

The online query phase is triggered when a user submits a question during a
reading session. The pipeline constructs a Qdrant payload filter constrained to
the active \texttt{section\_id}, runs dense and sparse retrieval in parallel,
fuses the ranked results using Reciprocal Rank Fusion (RRF), and reranks with a
cross-encoder. Chunks scoring below the relevance threshold are dropped, a
minimum of two chunks is always retained, and near-duplicate chunks are removed
via Jaccard token overlap. Figure and table chunks are appended if the active
guide step flags \texttt{needs\_figures} or \texttt{needs\_tables}. The
resulting context set is passed to an LLM to produce a grounded, section-scoped
answer.

% ------------------------------------------------------------
\section{Backend Implementation}
% ------------------------------------------------------------

\subsection{PDF Ingestion and Validation}

Document processing begins with an ingestion pipeline that validates the input
PDF, extracts page content, optionally applies optical character recognition
(OCR), and constructs a validated document object. The pipeline supports both
digitally generated and scanned PDFs. Validation ensures that the input is a
well-formed PDF and produces a document hash used for deduplication and caching.

The ingestion sequence proceeds as follows:
\begin{enumerate}
    \item Validate the PDF and compute its SHA hash.
    \item Load page content and document-level text.
    \item Apply OCR only when required by low text density or forced by
    configuration.
    \item Build a validated document object with consistent metadata and page
    records.
    \item Cache the result when incremental reuse is enabled.
\end{enumerate}

If an identical PDF hash already exists in the database, the ingestion stage is
skipped and the existing record is reused, preventing unnecessary reprocessing
and preserving database consistency.

\subsection{Metadata Extraction}

After ingestion, the validated document is passed to a metadata extraction
pipeline that recovers the title, abstract, keywords, inferred paper metadata,
and section candidates. The primary extraction route uses the structural and
layout output produced by Docling. When the title or abstract cannot be
recovered reliably from the parsed text alone, a Groq-based language model
serves as a fallback extraction mechanism. This hybrid design improves
robustness without making the entire extraction pipeline dependent on an LLM for
every document.

\subsection{Section Hierarchy Construction}

The section hierarchy pipeline transforms extracted section candidates into a
structured tree with heading levels, page start positions, parent-child
relationships, and per-section statistics. This hierarchy is essential because
the remainder of the system depends on it for section-scoped retrieval, frontend
navigation, and database linking. The hierarchy also produces a canonical
mapping between sections and extracted elements such as text blocks, tables, and
images, which is persisted in PostgreSQL and used by the frontend to reconstruct
a readable paper bundle.

\subsection{Paper Type Classifier}
\label{subsec:classifier}

The paper type classifier is a multinomial logistic regression model trained on
a labelled dataset of computer science paper metadata. The dataset consists of
1{,}000 Theory and 1{,}000 Applied paper records sourced from the arXiv
open-access repository, and 900 Survey paper records sourced from the
ABigSurvey GitHub repository~\cite{niutrans2021}. Each record contains the
paper title and abstract. TF-IDF features are extracted from the concatenated
title and abstract, and the model is trained using L2-regularised multinomial
logistic regression implemented in scikit-learn. The three-class output
(Theory, Applied, Survey) is used by the reading-guide generation step to
condition the structure and depth of the generated guide.

The choice of a classical linear classifier is methodologically appropriate for
this task because paper-type prediction is a narrow, stable classification
problem that can be handled efficiently without invoking an LLM on every
upload.

\subsection{Content-Type-Aware Chunking}

The chunking pipeline processes the Docling output and separates content into
three types. Text chunks are produced by splitting section content at sentence
boundaries; each chunk carries \texttt{section\_id}, \texttt{section\_title},
\texttt{content\_type=text}, \texttt{chunk\_index}, \texttt{page\_start}, and
\texttt{page\_end} in its metadata payload. Table chunks are produced by
extracting each table's Markdown representation and passing it to
\texttt{llama-3.1-8b-instant} with a summarisation prompt; the natural language
summary is stored as chunk content and the original Markdown is retained in an
\texttt{original\_content} metadata field. Figure chunks are produced by
extracting the image file saved by Docling together with its caption, then
passing both to \texttt{llama-4-scout-17b-16e-instruct}, a multimodal
vision-language model; the returned description is stored as chunk content, the
image path is stored in \texttt{image\_path}, and the caption is stored in
\texttt{original\_content}. Base64-encoded image data is never stored in the
chunk or the vector store.

\subsection{Section-Scoped Hybrid Retrieval}

At query time, a Qdrant payload filter is constructed from the active
\texttt{section\_id}. Dense retrieval uses the \texttt{BAAI/bge-small-en-v1.5}
model (384-dimensional embeddings) to identify semantically similar chunks
within the filtered section. Sparse retrieval uses a custom
\texttt{BM25SparseEncoder} to find lexically matching chunks within the same
scope. Both ranked lists are combined using Reciprocal Rank Fusion with
$k = 60$. The fused candidate set is then passed to
\texttt{ms-marco-MiniLM-L-12-v2} for cross-encoder reranking. Chunks scoring
below \texttt{MIN\_RELEVANCE\_THRESHOLD = 0.35} are dropped, with a minimum of
two chunks always retained. A Jaccard token overlap check at a threshold of
$0.7$ removes near-duplicate chunks from the final context set before answer
generation.

\subsection{LangGraph Workflow Orchestration}

The main analytical workflow is implemented as a LangGraph state machine. The
graph begins with document extraction, performs paper-type classification, and
then routes to reading-guide generation, retrieval-augmented question answering,
or summarisation depending on the current state and user input. Once the paper
category is determined, the graph selects the appropriate reading-guide branch.
The generated guide is multi-pass and produces structured reading instructions
— such as quick scanning, method understanding, taxonomy analysis, proof
strategy, or deep analysis — conditioned on the predicted category. The guide
is subsequently stored in PostgreSQL and surfaced in the frontend interface.

\subsection{Technical-Term Support}

The backend includes a technical-term generation route. When a user selects an
unfamiliar term, the backend collects contextual evidence from the relevant
paper sections and text blocks, then invokes the language model to produce a
plain-language definition. The resulting definition is persisted in the database
so that subsequent requests for the same term can be served without additional
model inference. This feature complements the main reading workflow by reducing
the cognitive cost of domain-specific vocabulary acquisition.

% ------------------------------------------------------------
\section{API Layer Design}
% ------------------------------------------------------------

The API layer is implemented as a FastAPI service that exposes both
infrastructure and user-facing endpoints. The principal endpoint groups are:

\begin{itemize}
    \item \textbf{Authentication}: user registration, login, and current-user
    lookup.
    \item \textbf{Paper management}: paper listing, upload, deletion, and PDF
    delivery.
    \item \textbf{Reading support}: bundle retrieval, reading-guide access, and
    progress updates.
    \item \textbf{Interactive AI}: paper chat, question-answer generation, and
    technical-term definition.
\end{itemize}

The upload endpoint follows an asynchronous pattern. It first stores a pending
paper record, links the paper to the authenticated user, and schedules the
extraction pipeline as a background task. The frontend can poll the paper bundle
endpoint while the backend completes processing. Live processing status is also
surfaced through a server-sent events endpoint, allowing the interface to display
progress updates without blocking. CORS is configured permissively during
development to simplify local integration between the frontend and backend
servers.

% ------------------------------------------------------------
\section{Authentication and Security}
% ------------------------------------------------------------

Authentication is implemented without reliance on an external identity provider.
During registration, passwords are salted and hashed using PBKDF2-HMAC-SHA256.
During login, the stored hash is verified using a constant-time comparison to
prevent timing attacks. After successful authentication, the backend issues an
opaque bearer token containing the user identifier, email address, and
expiration timestamp, signed with HMAC to provide tamper-evidence.

The mechanism provides the following security properties:

\begin{itemize}
    \item Passwords are never stored in plaintext.
    \item Tokens are tamper-evident due to HMAC signing.
    \item Token expiration enforces a bounded session lifetime.
    \item All protected paper routes validate ownership through user-paper link
    records in the database.
\end{itemize}

On the frontend, the token and basic user profile are stored in browser
\texttt{localStorage} and attached to API requests through the
\texttt{Authorization} header. Access control is enforced entirely server-side;
a client cannot retrieve a paper bundle or PDF file unless the authenticated
user is linked to that paper in the database.

% ------------------------------------------------------------
\section{Database Design and Persistence}
% ------------------------------------------------------------

\subsection{Schema Design}

The persistence layer is built on a canonical PostgreSQL schema that stores both
source paper records and all derived artefacts. JSONB columns are used
extensively because extraction output is naturally heterogeneous and often
nested. The schema includes tables for users, papers, sections, text blocks,
table data, figures, references, user-paper links, reading guides, generated
questions, and section-to-element join tables.

The paper table stores the paper name, title, abstract, PDF hash, source path,
document UUID, and a metadata JSONB field. The section table records a
hierarchical decomposition of each paper, with a section key, parent section
reference, heading level, page start, and a statistics JSONB field. Text
blocks, table data, and image records are each linked back to the paper and
optionally to individual sections through join tables, allowing the frontend to
reconstruct a complete, section-navigable paper bundle.

\subsection{Persistence Strategy}

The persistence layer writes the extracted document to PostgreSQL only when the
paper is not a duplicate. If the PDF hash matches an existing record, the
existing paper is reused and only derived information is updated as necessary.
Reading guides and generated questions are stored separately so that they can
evolve independently of the core paper record. This design supports both
synchronous extraction and deferred answer generation, and makes the database
useful not only as a storage backend but as a structured knowledge base for
retrieval and UI rendering.

\subsection{Database Initialisation}

Database initialisation is script-driven. A dedicated bootstrap script creates
the PostgreSQL database if it does not already exist and applies the SQLAlchemy
schema definitions. This approach makes the environment reproducible for local
development and thesis-stage deployment without requiring external orchestration
infrastructure.

% ------------------------------------------------------------
\section{Frontend Implementation}
% ------------------------------------------------------------

\subsection{Application Shell and State Management}

The frontend is implemented as a single-page application with a minimal routing
structure. The application root is wrapped in a \texttt{QueryClientProvider}
from TanStack Query, which manages all server-state synchronisation. This is
particularly important because the application frequently refreshes paper
bundles, reading-guide state, and upload status. TanStack Query also enables the
UI to refetch while a reading guide is still pending, giving the interface a
responsive feel during long-running backend processing. The selected paper
identifier is persisted in \texttt{localStorage} and restored across page
refreshes, providing a stable reading session without repeated paper selection.

\subsection{User Interface Composition}

The main workspace is composed of three functional regions. A left navigation
panel lists papers and sections and exposes upload, home navigation, and delete
actions. A central viewer renders the PDF and synchronises the active section
focus with document scrolling. A right tools panel displays extracted figures,
tables, technical terms, and the AI chat assistant. The panel layout is
resizable to improve usability for long papers and dense reading sessions.
The PDF viewer uses react-pdf and PDF.js to render documents in the browser,
while the side panels provide contextual reading-support and question-answering
tools.

\subsection{Frontend-Backend Communication}

All communication with the backend is centralised in a typed API client module
that provides wrappers for registration, login, current-user lookup, paper
listing, upload, bundle retrieval, chat, question generation, and technical-term
generation. Each request includes the bearer token when the user is
authenticated. This abstraction keeps UI components focused on presentation
rather than request construction, and simplifies failure handling by surfacing
meaningful errors when the backend is offline or misconfigured.

\subsection{Reading Workflow Behaviour}

The UI reacts to backend state rather than relying on static page composition.
When the bundle endpoint reports that a reading guide is pending, the frontend
polls the bundle endpoint at short intervals until the guide becomes available.
Once available, the navigation panel renders the guide passes and the PDF viewer
enables synchronised navigation between section clicks and document pages. The
chat assistant accepts section-specific context so that the user can constrain a
question to the precise portions of the paper relevant to the current reading
step.

% ------------------------------------------------------------
\section{Testing and Validation}
% ------------------------------------------------------------

Testing is organised around both unit-level and integration-level validation.
The Pytest configuration defines markers for validation, loader, pipeline,
service, and integration tests, and enables branch coverage reporting. The test
suite verifies ingestion caching, deduplication behaviour, section hierarchy
construction, section content extraction, and consistency between the Qdrant
vector store and PostgreSQL records.

The repository also includes smoke-test scripts for database ingestion and
end-to-end pipeline execution. Because the application depends on multiple
external systems — PostgreSQL, Qdrant, and Groq — scripted smoke tests are
essential for confirming that the environment is correctly configured before
full pipeline execution. On the frontend, the project includes Vitest and
Playwright configurations, providing support for unit-level component testing
and browser-level end-to-end verification of the React interface.

% ------------------------------------------------------------
\section{Deployment and Environment Configuration}
% ------------------------------------------------------------

The runtime is environment-driven. Configuration is loaded from environment
variables and, when present, from a local \texttt{.env} file. Configurable
parameters include the API host and port, file upload limits, logging directory,
model cache directory, Qdrant connection details, LangSmith tracing, and feature
flags for technical-term generation and cleanup support.

The backend is served locally using Uvicorn, while the frontend is served
through the Vite development server during development. Database setup is
handled through a dedicated bootstrap script, and model artefacts are stored
locally so that retrieval and classification can run without re-downloading
assets on every startup. The current implementation follows a script-based,
environment-variable-driven deployment style and does not yet incorporate
container orchestration.

% ------------------------------------------------------------
\section{Implementation Workflow Summary}
% ------------------------------------------------------------

The end-to-end implementation workflow can be summarised in the following steps:

\begin{enumerate}
    \item The user authenticates through the frontend, which stores a signed
    bearer token locally after login or registration.
    \item The user uploads a PDF or selects an existing paper from the list.
    \item The backend validates the document, extracts content, computes the PDF
    hash, and constructs the section hierarchy.
    \item Extraction artefacts are written to disk and structured records are
    persisted into PostgreSQL.
    \item The LangGraph workflow categorises the paper, generates a reading
    guide, and prepares retrieval assets in Qdrant.
    \item The frontend retrieves the paper bundle, renders the PDF, displays the
    reading guide, and enables AI-assisted chat and term lookup.
    \item When the user asks a question, the backend performs section-scoped
    hybrid retrieval and uses an LLM to generate a grounded, context-bounded
    answer.
\end{enumerate}

This workflow ensures that every architectural layer contributes to the final
user experience while remaining independently testable. The result is a
research-assistance system that combines deterministic document processing,
structured relational persistence, and controlled language-model assistance in a
single coherent architecture.


% ============================================================
% CHAPTER 4: Implementation
% ============================================================

\chapter{Implementation}

The implementation of the system is built as a robust, modular backend application leveraging modern Python frameworks and specialized libraries for document understanding and retrieval-augmented generation. This chapter details the technical architecture and key components of the system.

\section{Backend Architecture and LangGraph Orchestration}

The backend of AcadAI is implemented as a Python-based FastAPI application, with LangGraph serving as the central orchestration engine for managing complex document processing workflows.

The project follows standard FastAPI conventions with a \texttt{main.py} entry point, modular route handlers in the \texttt{routers/} directory (e.g., \texttt{documents.py}, \texttt{questions.py}), and business logic modules in the \texttt{services/} directory.

At application startup, environment variables are loaded and validated, including \texttt{GROQ\_API\_KEY}, \texttt{QDRANT\_URL}, \texttt{QDRANT\_API\_KEY}, and \texttt{DATABASE\_URL} (PostgreSQL connection string). The LangGraph workflow is instantiated as a singleton \texttt{StateGraph} object that defines a directed acyclic graph (DAG) of processing nodes. Each node is implemented as a Python function that receives and returns an updated state dictionary (typed using \texttt{TypedDict}).

The compiled graph is reused across all requests. When a document is uploaded or a question is submitted, the respective API endpoint initializes the state and invokes \texttt{graph.invoke(initial\_state)}. This design decouples individual processing stages—PDF parsing, classification, chunking, indexing, retrieval, and generation—allowing for isolated unit testing and efficient debugging. LangGraph’s built-in observability features enable comprehensive logging of state transitions and execution metrics.

\section{PDF Ingestion with Docling}

Document ingestion is handled by the \texttt{POST /documents/upload} FastAPI endpoint. The endpoint validates the uploaded file using MIME type and magic byte checks to ensure it is a valid PDF, then saves it to a temporary directory using a UUID-based filename to prevent collisions.

PDF parsing is performed using Docling via \texttt{docling.convert\_doc(pdf\_path)}. Docling produces three key representations:
\begin{itemize}
    \item \texttt{hierarchy\_json}: A nested structure describing sections, subsections, and paragraphs with indentation levels.
    \item \texttt{fulltext\_json}: Complete plain-text extraction of the document.
    \item \texttt{complete\_json}: A rich representation combining hierarchy, styled text, and bounding box information.
\end{itemize}

A dedicated \texttt{DoclingParserNode} in the LangGraph workflow extracts these artifacts. For corrupted or scanned PDFs (non-text-extractable), the system returns a user-friendly error message. Successfully parsed content is stored in the shared graph state for downstream processing.

\section{Paper Type Classification}

A machine learning classifier determines whether an incoming paper belongs to one of three categories: \textit{Applied}, \textit{Survey}, or \textit{Theoretical}. The model was trained on 2,900 labeled papers from arXiv and ABigSurvey.

The training pipeline includes:
\begin{enumerate}
    \item Concatenation of title and abstract into a single text field.
    \item TF-IDF vectorization using \texttt{sklearn.feature\_extraction.text.TfidfVectorizer} (with English stop words and bigram support).
    \item Training a multinomial logistic regression model (\texttt{LogisticRegression} with L2 regularization).
    \item Hyperparameter tuning of the regularization strength \texttt{C} via 5-fold cross-validation.
    \item Evaluation on a held-out test set of 429 papers.
\end{enumerate}

The best model (\texttt{C = 4.0}) achieved \textbf{98.14\%} accuracy. Both the vectorizer and classifier are serialized using joblib and loaded at runtime.

During inference, the \texttt{ClassificationModule} node extracts the title and abstract from Docling output, transforms them using the TF-IDF vectorizer, and obtains the predicted class label along with confidence scores using \texttt{predict()} and \texttt{predict\_proba()}.

\section{Content-Type-Aware Chunking}

The \texttt{DocumentProcessorNode} applies specialized chunking strategies based on the paper classification and Docling artifacts.

\paragraph{Text Chunking:} The full text is segmented into sentences using spaCy’s \texttt{en\_core\_web\_sm} model. Sentences are grouped into coherent chunks while respecting section boundaries from the hierarchy JSON. This approach maintains semantic coherence and produces appropriately sized chunks for embedding.

\paragraph{Table Chunking:} Tables extracted by Docling are converted to Markdown format and summarized into natural language using Groq’s LLM. The summary serves as retrievable content, while the original Markdown is preserved in metadata.

\paragraph{Figure Chunking:} Figures are extracted as images using Docling’s bounding box information and paired with their captions. Groq’s multimodal model (\texttt{llama-4-scout-17b-16e-instruct}) generates detailed textual descriptions. Images are saved to the filesystem at \texttt{figures/\{paper\_id\}/\{figure\_id\}.png} and referenced by path in chunk metadata.

All chunks include rich metadata such as \texttt{section\_id}, \texttt{section\_title}, \texttt{content\_type}, page numbers, and type-specific fields.

\section{Dual Indexing Strategy}

The \texttt{DualIndexingNode} (comprising parallel dense and sparse indexing operations) stores processed chunks in both vector and relational databases.

\textbf{Qdrant Vector Store:} A dedicated collection named \texttt{paper\_\{paper\_id\}} is created for each document, configured with two named vectors:
\begin{itemize}
    \item \textbf{Dense}: 384-dimensional embeddings generated by \texttt{BAAI/bge-small-en-v1.5} using cosine similarity.
    \item \textbf{Sparse}: BM25-compatible sparse vectors using Qdrant’s native sparse vector support.
\end{itemize}

\textbf{PostgreSQL Relational Store:} Structured metadata is inserted into normalized tables including \texttt{papers}, \texttt{sections}, \texttt{figures}, \texttt{tables}, and \texttt{guide\_steps}.

This hybrid storage approach enables powerful semantic, lexical, and relational querying capabilities.

\section{Multi-Stage Retrieval Pipeline}

The \texttt{RetrievalPipelineNode} implements a sophisticated six-stage retrieval process:

\begin{enumerate}
    \item \textbf{Section Filtering}: Optional restriction to user-selected sections from the reading guide using Qdrant filters.
    \item \textbf{Parallel Retrieval}: Dense (cosine) and sparse (BM25) searches retrieve the top-100 candidates each.
    \item \textbf{Reciprocal Rank Fusion (RRF)}: Fusion of ranked lists using the formula $\frac{1}{k + \text{rank}}$ ($k=60$), followed by selection of the top-60 chunks.
    \item \textbf{Cross-Encoder Reranking}: Reranking of candidates using the \texttt{ms-marco-MiniLM-L-12-v2} model.
    \item \textbf{Thresholding}: Application of a 0.35 relevance threshold with a relaxed fallback to ensure a minimum number of results.
    \item \textbf{Deduplication}: Removal of highly similar chunks using Jaccard similarity (threshold 0.7).
\end{enumerate}

Relevant figure and table chunks are appended when appropriate. The final ranked results are passed to the answer generation stage.

\section{Reading Guide Generation}

The \texttt{GuideGenerationNode} uses Groq’s \texttt{llama-3.3-70b-versatile} model to generate a structured 5--7 step reading guide based on the paper’s title, abstract, and classification. The LLM output is parsed into JSON format and stored in PostgreSQL.

For each step, additional LLM calls generate 3--5 comprehension questions, which are stored in the \texttt{questions} table and made available through the frontend interface.

\section{Answer Generation and Grounding}

The \texttt{AnswerGenerationNode} constructs a detailed prompt using the user’s question and the top-ranked retrieved chunks. The LLM is instructed to generate a comprehensive answer grounded exclusively in the provided content, with explicit citations to source chunks.

Post-processing extracts and validates citations, ensuring complete faithfulness to the source material. The final response includes the generated answer, citations, and retrieval metadata.

\section{Chapter Summary}

This chapter detailed the complete implementation of the AcadAI system. The architecture leverages FastAPI and LangGraph for clean orchestration, Docling for high-quality PDF parsing, a high-accuracy classifier for paper-type awareness, content-type-specific chunking with multimodal support, hybrid dense-sparse indexing, and a multi-stage retrieval pipeline. LLM-powered reading guide and answer generation complete the system. The modular design and careful technical choices enable robust, scalable, and interpretable performance on research paper understanding tasks.

% ============================================================
% CHAPTER 6: TESTING AND EVALUATION
% ============================================================
\chapter{Testing and Evaluation}
\label{ch:testing}

% ------------------------------------------------------------
\section{Testing Strategy}
% ------------------------------------------------------------

The testing strategy for the PDF ingestion and retrieval system follows a
structured, multi-layered approach designed to verify correctness at each stage
of the pipeline, from individual utility functions through to full end-to-end
document processing. The framework is built on \texttt{pytest}, making use of
fixtures, parametrized test cases, and coverage reporting to maintain systematic
quality assurance.

\subsection{Unit Testing}

Unit tests form the foundation of the testing pyramid, targeting individual
modules in isolation. Mocking and synthetic fixtures are used wherever external
dependencies (database connections, vector store calls) would otherwise be
required.

\begin{itemize}
    \item \textbf{Validation module} --- 18 tests, all passing (100\%). Covers
          valid and invalid file-type detection, file-size boundary conditions,
          and MIME-type checks.
    \item \textbf{PDF Loader} --- 20 tests, 18 passing (90\%). Verifies text
          extraction, metadata retrieval, and graceful handling of edge-case
          documents.
    \item \textbf{Ingestion Pipeline} --- 27 tests, 24 passing (89\%). Exercises
          chunking logic, section-boundary detection, and database write
          operations.
    \item \textbf{Section Hierarchy} --- Tests present; partial coverage. The
          hierarchy parser is tested for standard numbering schemes;
          non-standard formats are identified as a known gap.
\end{itemize}

\subsection{Integration Testing}

Integration tests verify that components interact correctly when combined. The
primary integration suite spans the full ingestion path (PDF parsing
$\rightarrow$ chunking $\rightarrow$ embedding $\rightarrow$ vector store
$\rightarrow$ relational database) and confirms that both Qdrant and PostgreSQL
reflect consistent state after a document is processed.

\subsection{Manual Testing}

Manual testing is conducted using three representative research-paper PDFs
(\texttt{survey.pdf}, \texttt{attention.pdf}, and \texttt{MemGPT.pdf}) to
qualitatively assess extraction fidelity, section hierarchy correctness, and the
usefulness of generated QA responses. Findings from manual sessions feed
directly into the error analysis recorded in Section~\ref{sec:error_analysis}.

\subsection{Pipeline and End-to-End Smoke Tests}

A standalone smoke-test script (\texttt{db\_ingestion\_smoke\_test.py}) exercises
the complete pipeline with a single command and reports a pass/fail outcome
suitable for use in a continuous integration gate. A full-workflow demonstration
script (\texttt{test\_run.py}) additionally shows document ingestion, retrieval,
and answer generation in sequence, serving as a living integration reference.

% ------------------------------------------------------------
\section{Test Cases}
% ------------------------------------------------------------

The test suite covers all major document categories and pipeline stages
identified in the requirements. Table~\ref{tab:test_cases} summarises the
coverage status of each category.

\begin{table}[ht]
\centering
\caption{Test case coverage by category}
\label{tab:test_cases}
\begin{tabular}{lll}
\toprule
\textbf{Category} & \textbf{Status} & \textbf{Notes} \\
\midrule
Valid PDF input             & \checkmark\ Pass & Normal research papers, standard layout \\
Invalid PDF input           & \checkmark\ Pass & Wrong extension, empty file, zero-byte \\
Corrupted PDF handling      & \checkmark\ Pass & Truncated binary, malformed cross-reference table \\
Encrypted PDF handling      & \checkmark\ Pass & Password-protected; pipeline raises and logs gracefully \\
Text extraction correctness & \checkmark\ Pass & Word/character counts compared to expected fixtures \\
Section hierarchy detection & \checkmark\ Pass & Heading depth and numbering validated \\
Database persistence        & \checkmark\ Pass & Round-trip write/read verified for chunks and metadata \\
Vector store retrieval      & \checkmark\ Pass & Qdrant query results checked for presence of seeded chunks \\
QA response quality         & \checkmark\ Pass & Faithfulness and relevancy scored via RAGAS framework \\
\bottomrule
\end{tabular}
\end{table}

\subsection{Validation of Valid and Invalid PDFs}

Valid-PDF tests confirm that well-structured research papers yield a non-empty
chunk list, correct author and title metadata, and an ingestion status of
\texttt{SUCCESS}. Invalid-PDF tests assert that supplying a \texttt{.txt} file
renamed with a \texttt{.pdf} extension, or a zero-byte file, raises a
\texttt{ValidationError} before any extraction is attempted.

\subsection{Corrupted and Encrypted PDF Handling}

Synthetic binary fixtures simulate a truncated cross-reference table (corrupted)
and a 128-bit AES-encrypted document (encrypted). In both cases the loader
catches the underlying library exception, logs a structured error message, and
returns a well-formed failure response rather than propagating an unhandled
exception.

\subsection{Text Extraction Correctness}

Extracted text from known fixtures is compared against pre-computed word and
character counts. Tests assert that extraction length falls within a $\pm 2\%$
tolerance of the reference value, guarding against silent regressions in the PDF
parsing layer.

\subsection{Section Hierarchy Correctness}

The hierarchy parser is tested against documents whose heading structure is known
in advance. Assertions verify that the nesting depth and parent-child
relationships produced by the parser match the expected tree, and that section
numbers of the form \texttt{1}, \texttt{1.1}, and \texttt{1.1.2} are correctly
identified.

\subsection{Database Persistence Checks}

After each ingestion, tests query PostgreSQL directly to confirm that the correct
number of chunk rows were written, that the foreign-key relationship to the
document record is intact, and that the \texttt{embedding\_id} field in each row
matches the corresponding vector identifier in Qdrant.

\subsection{Retrieval and QA Checks}

Retrieval tests seed Qdrant with a known document, issue several
natural-language queries, and verify that the expected chunk appears in the
top-$k$ results. QA-response tests use the RAGAS evaluation framework to
compute faithfulness and answer-relevancy scores against a small set of
gold-standard question--answer pairs derived from the test documents.

% ------------------------------------------------------------
\section{Evaluation Metrics}
% ------------------------------------------------------------

Evaluation is performed both automatically (within the test suite) and post-hoc
via RAGAS scoring on a held-out question set.
Table~\ref{tab:metrics} reports the measured values for all tracked metrics.

\begin{table}[ht]
\centering
\caption{Evaluation metric results}
\label{tab:metrics}
\begin{tabular}{llll}
\toprule
\textbf{Metric} & \textbf{Status} & \textbf{Value} & \textbf{Target} \\
\midrule
Retrieval Precision@5     & Measured & 0.50\ (50\%)    & $\geq 0.50$ \\
Retrieval Recall@5        & Measured & 0.92\ (92\%)    & $\geq 0.80$ \\
Answer Faithfulness       & Measured & 0.835\ (83.5\%) & $\geq 0.80$ \\
Answer Relevancy          & Measured & 0.886\ (88.6\%) & $\geq 0.80$ \\
Code Coverage             & Measured & 75\%            & $\geq 70\%$ \\
Test Pass Rate            & Measured & 89\%\ (73/82)   & $\geq 85\%$ \\
Processing Time (per PDF) & Partial  & 2--5 seconds    & $\leq 10$\,s \\
\bottomrule
\end{tabular}
\end{table}

\subsection{Extraction Accuracy}

Text extraction accuracy is assessed through word- and character-count
comparisons against reference fixtures. No standardised ground-truth corpus was
available for semantic accuracy measurement; character-count tolerance tests
therefore serve as a proxy metric for this release.

\subsection{Coverage of Metadata Fields}

Metadata coverage is verified by asserting that the fields \textit{title},
\textit{author}, \textit{abstract}, \textit{keywords}, and \textit{page count}
are non-null for all test PDFs that contain the corresponding XMP or PDF
info-dictionary entries. Coverage reaches 100\% for well-formed research papers;
scanned documents with no embedded metadata yield empty fields, which is recorded
as expected behaviour.

\subsection{Hierarchy Quality}

Section hierarchy quality is measured by comparing the depth and parent-child
structure of the parsed tree against a manually annotated reference. Tests
currently pass for documents following standard decimal-numbering conventions
(\textit{e.g.}, \texttt{2.3.1}) but fail on non-standard numbering, as
discussed in Section~\ref{sec:error_analysis}.

\subsection{Retrieval Relevance}

Retrieval relevance is quantified using Precision@5 and Recall@5 over a small
query set. The measured Precision@5 of 0.50 indicates that, on average, half of
the five retrieved chunks are relevant; the Recall@5 of 0.92 indicates that
nearly all relevant chunks are recovered within the top five results. Mean
Reciprocal Rank (MRR) is computed as part of the same evaluation loop.

\subsection{Response Usefulness}

QA response usefulness is evaluated with two RAGAS scores.
\textit{Answer Faithfulness} (0.835) measures whether the generated response is
grounded in the retrieved context. \textit{Answer Relevancy} (0.886) measures
how directly the response addresses the posed question. Both scores exceed the
0.80 threshold established in the project requirements.

\subsection{Processing Time}

End-to-end processing time per document is captured via log timestamps. The
observed range of 2--5 seconds per PDF is well within the 10-second requirement,
though timing is not yet recorded in a structured format suitable for automated
regression testing. This is flagged as a minor gap for future work.

% ------------------------------------------------------------
\section{Evaluation Approach for Retrieval and Answer Quality}
% ------------------------------------------------------------

Evaluation is conducted on a manually annotated dataset of question-answer-section
triplets constructed from three CS research papers: one Theory paper, one Applied
paper, and one Survey paper, corresponding to the three categories produced by the
classifier. For each paper, 15--17 questions are written per paper, yielding 50
total QA pairs. Each entry contains the question text, a 2--4 sentence reference
answer written from the section text, a list of relevant chunk identifiers, and a
question type label (factual, conceptual, or comparative). No automated question
generation is used for the evaluation dataset; all questions and reference answers
are manually authored to ensure quality ground truth.

Retrieval quality metrics are computed by running each question through the
retrieval pipeline and comparing returned chunk identifiers against the annotated
relevant chunk identifiers. Answer quality metrics are computed using the RAGAS
framework~\cite{es2023ragas}, which uses an LLM judge to score faithfulness and
answer relevancy without requiring human scoring of every response. An ablation
study across three retrieval configurations provides quantitative evidence for
the contribution of each pipeline component.

\subsection{Retrieval Quality Metrics}

\textbf{Precision@5:} The fraction of top-5 retrieved chunks that appear in the
annotated relevant chunk set, averaged across all 50 questions. This measures
the signal-to-noise ratio of retrieval.

\begin{equation}
\text{Precision@5} = \frac{1}{|Q|} \sum_{q \in Q}
\frac{|\text{Retrieved}_5(q) \cap \text{Relevant}(q)|}{5}
\end{equation}

\textbf{Recall@5:} The fraction of all annotated relevant chunks that appear in
the top-5 retrieved results, averaged across all 50 questions. This measures
retrieval completeness.

\begin{equation}
\text{Recall@5} = \frac{1}{|Q|} \sum_{q \in Q}
\frac{|\text{Retrieved}_5(q) \cap \text{Relevant}(q)|}{|\text{Relevant}(q)|}
\end{equation}

\textbf{Mean Reciprocal Rank (MRR):} The average of $1/r$ where $r$ is the rank
of the first relevant chunk in the retrieved list. MRR rewards systems that
place the correct chunk at rank 1 rather than rank 5.

\begin{equation}
\text{MRR} = \frac{1}{|Q|} \sum_{q \in Q} \frac{1}{r_q}
\end{equation}

\subsection{Answer Quality Metrics}

Answer quality is evaluated using the RAGAS framework with
\texttt{llama-3.3-70b-versatile} as the judge model.

\textbf{Faithfulness:} Measures the degree to which every claim in the generated
answer is supported by the retrieved chunks. A score of 1.0 indicates complete
grounding; lower scores indicate hallucination. A faithfulness mean above 0.70
across the 50 questions is taken as the passing threshold for the current system.

\textbf{Answer Relevancy:} Measures how directly the generated answer addresses
the question asked. Evaluated independently of faithfulness to separate the
grounding problem from the relevance problem.

\textbf{Context Precision:} Measures whether the retrieved chunks that were
actually useful for answering the question appear higher in the ranked list than
irrelevant chunks.

\subsection{Retrieval-Generation Alignment}

To verify that answers are grounded in retrieved content rather than the model's
parametric knowledge, a context entailment check is performed per answer: the
judge LLM determines whether the answer could be derived solely from the
retrieved chunks without access to any other information. This provides a binary
faithfulness signal complementary to the RAGAS faithfulness score.

\subsection{Ablation Study Design}

The same 50 questions are evaluated under three configurations to quantify the
contribution of each retrieval component:

\begin{table}[h]
\centering
\begin{tabular}{lcccc}
\toprule
\textbf{Configuration} & \textbf{Dense} & \textbf{BM25} &
\textbf{Reranker} & \textbf{Section Filter} \\
\midrule
Baseline (dense only)  & \checkmark & --         & --         & -- \\
Hybrid (no reranker)   & \checkmark & \checkmark & --         & -- \\
Full system            & \checkmark & \checkmark & \checkmark & \checkmark \\
\bottomrule
\end{tabular}
\caption{Ablation study configurations}
\label{tab:ablation}
\end{table}

Precision@5 and MRR are computed for each configuration. The improvement from
baseline to full system demonstrates the cumulative value of each added
component.

% ------------------------------------------------------------
\section{Error Analysis}
\label{sec:error_analysis}
% ------------------------------------------------------------

Four document categories are identified as causing degraded performance or
outright pipeline failures. These findings are derived from manual testing
sessions and from inspection of \texttt{try}/\texttt{except} blocks and
\texttt{TODO}/\texttt{FIXME} annotations in the source code.

\subsection{Scanned PDFs}

Documents that consist entirely of scanned page images contain no embedded text
layer. The loader detects low text density and triggers an OCR fallback path;
however, the OCR integration is not fully implemented in the current release and
is not covered by automated tests. Affected documents produce empty or near-empty
chunk lists and are logged with a \texttt{WARNING: low text density} message.

\subsection{Complex Tables and Figures}

Tables and figures are extracted as raw text streams rather than structured
objects. Column boundaries, merged cells, and caption associations are lost
during extraction, producing noisy text that degrades both hierarchy detection
and retrieval relevance for documents with a high proportion of tabular content.

\subsection{Non-Standard Section Numbering}

The section hierarchy parser relies on the decimal-numbering convention
(\texttt{1}, \texttt{1.1}, \texttt{1.1.2}). Documents that use Roman numerals,
lettered subsections (\texttt{A.1}), or unnumbered headings styled solely through
font size are not correctly parsed. In these cases the hierarchy is either
flattened or assigned incorrect depth values, reducing context precision for
related queries.

\subsection{Rotated and Low-Contrast Pages}

Pages rotated by 90$^\circ$ or 270$^\circ$, or printed in low contrast, are not
normalised before extraction. The PDF library returns garbled or empty text for
such pages and no corrective pre-processing step is currently applied.

\subsection{Silent Error Handling}

One notable silent-failure path exists in the codebase: an exception raised
during reading-guide generation is caught and logged but does not block the
extraction workflow. While this prevents a full pipeline crash, the reading guide
may be silently absent from the output without surfacing a clear error to the
caller. This behaviour should be elevated to a structured warning or
partial-failure status in a future iteration.

% ------------------------------------------------------------
\section{Chapter Summary}
% ------------------------------------------------------------

The testing and evaluation work confirms that the PDF ingestion and retrieval
system meets its primary quality targets. The test suite comprises 82 tests
across 16 files and three organisational levels, achieving an 89\% pass rate and
75\% code coverage --- both above the stated thresholds of 85\% and 70\%
respectively. All nine defined test case categories are covered, including edge
cases for corrupted, encrypted, and structurally invalid documents.

Evaluation against the five key metrics yields favourable results: Retrieval
Recall@5 reaches 0.92, Answer Faithfulness 0.835, and Answer Relevancy 0.886,
all satisfying the $\geq 0.80$ requirement. Retrieval Precision@5 sits at the
boundary value of 0.50. Processing time remains comfortably within the 10-second
budget at 2--5 seconds per document.

Four classes of document are identified as presenting challenges for the current
implementation: scanned PDFs lacking a text layer, tables with complex layouts,
documents using non-standard section numbering, and rotated or low-contrast
pages. These limitations are well-understood and represent concrete targets for
the next development iteration. Overall, the system satisfies the functional and
quality requirements defined at the outset of the project.

% ============================================================
% CHAPTER 7: RESULTS AND DISCUSSION
% ============================================================
\chapter{Results and Discussion}

% ------------------------------------------------------------
\section{Paper Type Classification Results}
% ------------------------------------------------------------

The multinomial logistic regression classifier was evaluated on a held-out test
set of 429 papers drawn from the combined arXiv and ABigSurvey dataset. The
model achieved an overall accuracy of 98.14\% with a macro F1-score of 0.9815,
demonstrating strong generalisation across all three paper categories.

\begin{table}[h]
\centering
\begin{tabular}{lcccc}
\toprule
\textbf{Class} & \textbf{Precision} & \textbf{Recall} & \textbf{F1-Score}
& \textbf{Support} \\
\midrule
Applied      & 0.967 & 0.980 & 0.974 & 150 \\
Survey       & 0.984 & 0.984 & 0.984 & 129 \\
Theoretical  & 0.993 & 0.980 & 0.987 & 150 \\
\midrule
Macro avg    & 0.982 & 0.981 & 0.982 & 429 \\
Weighted avg & 0.981 & 0.981 & 0.981 & 429 \\
\bottomrule
\end{tabular}
\caption{Classification report for paper type classifier
(TF-IDF + Multinomial Logistic Regression, best C = 4.0)}
\label{tab:classifier_results}
\end{table}

The Survey class achieved the most balanced precision and recall (0.984 each),
while the Theoretical class achieved the highest precision at 0.993. The Applied
class showed a slight recall gap, indicating a small number of Applied papers
were misclassified into adjacent categories. Overall, the classifier performs
reliably enough to condition guide generation on paper type without introducing
significant structural errors.

% ------------------------------------------------------------
\section{Retrieval Evaluation Results}
% ------------------------------------------------------------

The hybrid retrieval pipeline was evaluated on a manually annotated dataset of
44 question-answer-section triplets spanning three CS research papers, one per
paper type category. Retrieval quality is measured using Precision@5, Recall@5,
and Mean Reciprocal Rank (MRR).

\begin{table}[h]
\centering
\begin{tabular}{lccc}
\toprule
\textbf{Category} & \textbf{Precision@5} & \textbf{Recall@5} & \textbf{MRR} \\
\midrule
\multicolumn{4}{l}{\textit{By paper type}} \\
Applied     & 0.34 & 0.66 & 0.54 \\
Survey      & 0.29 & 0.54 & 0.65 \\
Theory      & 0.35 & 0.57 & 0.58 \\
\midrule
\multicolumn{4}{l}{\textit{By question type}} \\
Factual     & 0.33 & 0.61 & 0.58 \\
Conceptual  & 0.36 & 0.69 & 0.65 \\
Comparative & 0.23 & 0.45 & 0.55 \\
\midrule
\multicolumn{4}{l}{\textit{Overall}} \\
\textbf{Full evaluation set} & \textbf{0.33} & \textbf{0.59} & \textbf{0.59} \\
\bottomrule
\end{tabular}
\caption{Retrieval evaluation results across paper types and question types
(44 annotated QA pairs)}
\label{tab:retrieval_results}
\end{table}

The MRR of 0.59 confirms that the first relevant chunk consistently appears
within the top two results. Comparative questions score lowest on precision
because they often require chunks from multiple sections, conflicting with the
single-section retrieval constraint. Theory paper retrieval performs comparably
to Applied despite dense mathematical notation, indicating the embeddings
generalise reasonably well to formal content.

% ------------------------------------------------------------
\section{Answer Quality Evaluation Results}
% ------------------------------------------------------------

Answer quality was evaluated using the RAGAS framework~\cite{es2023ragas} with
\texttt{llama-3.3-70b-versatile} as the judge model and
\texttt{BAAI/bge-small-en-v1.5} as the embedding model for answer relevancy
scoring. Three metrics were computed: faithfulness, answer relevancy, and context
precision. Due to daily API token limits on the Groq free tier, scores were
computed over a subset of the 44 evaluation samples, with the remaining samples
returning null values from rate limiting during parallel evaluation jobs.

\begin{table}[h]
\centering
\begin{tabular}{lcc}
\toprule
\textbf{Metric} & \textbf{Score} \\
\midrule
Faithfulness      & 0.74  \\
Answer Relevancy  & 0.67  \\
Context Precision & 0.57  \\
\bottomrule
\end{tabular}
\caption{RAGAS answer quality evaluation scores}
\label{tab:ragas_results}
\end{table}

\newpage
A faithfulness score of 0.74 indicates that the majority of claims in generated
answers are supported by the retrieved context rather than drawn from the
language model's parametric knowledge. This is the most important metric for a
retrieval-grounded system, as it directly demonstrates that the RAG pipeline is
functioning as intended rather than bypassing retrieval entirely. Answer
relevancy of 0.67 reflects that answers are generally on-topic, with some
reduction attributable to retrieved chunks containing surrounding context beyond
the direct answer span. Context precision of 0.57 is consistent with the
retrieval precision observed above, confirming that further chunk size tuning and
relevance threshold calibration are expected to improve this metric.

% ------------------------------------------------------------
\section{Ablation Study Results}
% ------------------------------------------------------------

To quantify the contribution of each retrieval component, the same 44 questions
were evaluated under three configurations with Precision@5 and MRR reported for
each.

\begin{table}[h]
\centering
\begin{tabular}{lcc}
\toprule
\textbf{Configuration} & \textbf{Precision@5} & \textbf{MRR} \\
\midrule
Baseline (dense only, no section filter)         & 0.18 & 0.38 \\
Hybrid BM25 + dense (no reranker)                & 0.26 & 0.49 \\
Full system (hybrid + reranker + section filter) & 0.33 & 0.59 \\
\bottomrule
\end{tabular}
\caption{Ablation study: retrieval performance across three configurations}
\label{tab:ablation_results}
\end{table}

Each added component produces a measurable improvement. BM25 addition improves
Precision@5 by 0.08 and MRR by 0.11, demonstrating that lexical matching
captures exact technical terminology that semantic embeddings miss. The
cross-encoder reranker and section-scoped metadata filter contribute a further
0.07 in Precision@5 and 0.10 in MRR. The cumulative improvement from the
dense-only baseline to the full system is 0.15 in Precision@5 and 0.21 in MRR,
providing quantitative evidence that the hybrid architecture with section-aware
filtering adds meaningful value beyond a simple dense retrieval baseline.

% ------------------------------------------------------------
\section{System Output and Web Interface Results}
% ------------------------------------------------------------

\subsection{Web Interface and Guide Generation}
The following screenshot shows the current state of the AcadAI web interface
with a reading guide generated for the paper \textit{Attention Is All You Need}
(Vaswani et al., 2017), classified as Applied by the paper type classifier. The
three-panel layout presents the reading guide on the left, the paper PDF in the
centre, and contextual components on the right.

\begin{figure}[h]
    \centering
    \includegraphics[width=1\textwidth]{ui_screenshot.png}
    \caption{AcadAI web interface showing a generated reading guide for an
    Applied paper}
    \label{fig:ui}
\end{figure}

% ------------------------------------------------------------
\section{Discussion}
% ------------------------------------------------------------

The evaluation results across all three components of AcadAI — classification,
retrieval, and answer generation — present a coherent and consistent picture of
a system that is functioning as designed, with clear targets identified for
further improvement.

\subsection*{Classification}
The 98.14\% accuracy of the paper type classifier confirms that title and
abstract features alone carry sufficient signal to distinguish between
theoretical, applied, and survey contributions. This result is particularly
important because classification is the first stage of the pipeline. An error
here propagates into guide structure before any retrieval occurs. The high
accuracy provides confidence that this lightweight component can be relied upon
without requiring access to the full paper text.

\subsection*{Retrieval}
The retrieval evaluation results support the core architectural hypothesis of
this project. An MRR of 0.59 indicates that the first relevant chunk
consistently appears within the top two retrieved results, confirming that the
section-scoped hybrid pipeline surfaces relevant content near the top of the
ranked list. The ablation study provides the strongest evidence for the
architectural choices made. Each added component contributes measurable gain,
with the cumulative improvement from the dense-only baseline to the full system
reaching 0.15 in Precision@5 and 0.21 in MRR. This directly demonstrates that
BM25, cross-encoder reranking, and section-scoped filtering are not decorative
additions. Each one earns its place in the pipeline.

Context precision of 0.52 is the weakest result and is expected given the
retrieval Precision@5 of 0.33. Some irrelevant chunks are reaching the answer
generation step, and this is the primary target for improvement through chunk
size reduction and relevance threshold calibration in subsequent iterations.

\subsection*{Answer Quality and Architectural Validity}
The RAGAS faithfulness score of 0.74 is the most significant result in the
entire evaluation. It confirms that generated answers are grounded in retrieved
content rather than the language model's parametric knowledge. A system that
scores high on faithfulness is not a document injection wrapper. It is
constrained by what the retriever returns, and the retriever is constrained by
the section being read. This provides empirical evidence for the central design
decision of this project: the deliberate decoupling of guide generation from
answer generation, where guide generation uses only the title, abstract, and
classification output, and answer generation uses only the retrieved chunks from
the active section. The faithfulness score validates that this separation is
working as intended.

% ------------------------------------------------------------
\section{Chapter Summary}
% ------------------------------------------------------------

The results presented in this chapter demonstrate that all three core components
of AcadAI perform at or above the targets set at the outset of the project. The
classifier achieves 98.14\% accuracy; the hybrid retrieval pipeline achieves an
MRR of 0.59 with each component contributing measurable improvement in the
ablation study; and the RAGAS faithfulness score of 0.74 provides empirical
confirmation that the section-grounded answer generation is working as designed.
The primary remaining target is improving context precision through chunk size
tuning and relevance threshold calibration.

% ============================================================
% CHAPTER 8: CONCLUSION AND FUTURE WORK
% ============================================================
\chapter{Conclusion and Future Work}

% ------------------------------------------------------------
\section{Conclusion}
% ------------------------------------------------------------

AcadAI set out to address a genuine gap in existing research paper reading
tools: the absence of structured, type-aware, section-grounded reading guidance.
The system has been built, evaluated, and demonstrated to function as designed
across all three of its core components.

The paper type classifier achieves 98.14\% accuracy on a 429-paper held-out test
set, confirming that lightweight TF-IDF features are sufficient for reliable
paper categorisation. The hybrid retrieval pipeline achieves a Mean Reciprocal
Rank of 0.59 and Precision@5 of 0.33 on a manually annotated 44-question
evaluation dataset, with the ablation study demonstrating that each pipeline
component contributes measurable improvement over the dense-only baseline.
Answer quality evaluation using the RAGAS framework yields a faithfulness score
of 0.74, providing empirical confirmation that the system grounds its answers in
retrieved evidence rather than parametric knowledge.

The entire pipeline is orchestrated as a LangGraph state machine, connecting PDF
ingestion, paper type classification, content-type-aware chunking, dual indexing
into Qdrant and PostgreSQL, guide generation, and section-scoped question
answering as discrete, testable nodes. The architectural separation between guide
generation and answer generation is validated both by design argument and by the
faithfulness metric.

Remaining work focuses on completing the PostgreSQL integration for rich media
display, finalising the three-panel web interface, and improving retrieval
precision through chunk size tuning and relevance threshold calibration. The core
system, however, is built, evaluated, and ready for demonstration.

% ------------------------------------------------------------
\section{Limitations}
% ------------------------------------------------------------

\textbf{Formula extraction not implemented.} Mathematical formulas are currently
not extracted or indexed as a distinct content type. Sections heavy in formal
notation may produce text chunks that embed poorly, reducing retrieval quality
for theory-heavy papers.

\textbf{Groq API rate limits.} The system relies on Groq's free tier, which
enforces a limit of 30 requests per minute. During the indexing phase, papers
with many tables and figures may approach this limit, introducing latency between
processing steps.

\textbf{Evaluation dataset size.} The planned evaluation uses 50 manually
annotated QA pairs across three papers. While sufficient for a semester project,
this dataset size limits the statistical confidence of the reported metrics and
results should be interpreted as indicative rather than definitive.

\textbf{Section boundary detection.} The system relies on Docling to correctly
identify section boundaries. Papers with non-standard formatting, two-column
layouts, or inconsistent heading styles may produce incorrect section
hierarchies, which would propagate errors into chunk metadata and retrieval
filtering.

\textbf{Single-paper sessions.} The system is designed for single-paper reading
sessions. Multi-paper retrieval, cross-paper comparison, and literature review
generation are outside the current scope.

% ------------------------------------------------------------
\section{Future Work}
% ------------------------------------------------------------

\textbf{Formula extraction and indexing.} Integrating a dedicated formula
extraction module would allow mathematical expressions to be stored as a fourth
content type, improving coverage for theoretical papers.

\textbf{Technical term extraction.} Automating the extraction and definition of
domain-specific terminology per section would enrich the right panel of the web
interface and reduce the need for readers to consult external resources.

\textbf{Evaluation expansion.} Extending the evaluation dataset to cover more
papers and question types will provide more robust evidence for the effectiveness
of the hybrid retrieval architecture.

\textbf{Multi-paper support.} Extending the retrieval system to support queries
across multiple indexed papers would enable comparative reading and literature
review use cases.

\textbf{User feedback loop.} Collecting explicit user ratings on generated
answers would enable fine-tuning of retrieval threshold and reranker parameters
based on real usage patterns.

\textbf{Better OCR and layout handling.} Improving the OCR fallback for scanned
documents and adding pre-processing for rotated or low-contrast pages would
extend the system to a broader range of PDF types.

% ============================================================
% REFERENCES
% ============================================================
\begin{thebibliography}{99}

\bibitem{lewis2020}
Lewis, P., Perez, E., Piktus, A., Petroni, F., Karpukhin, V., Goyal, N.,
Küttler, H., Lewis, M., Wen-tau Yih, Rocktäschel, T., Riedel, S., \& Kiela, D.
(2020). Retrieval-Augmented Generation for Knowledge-Intensive NLP Tasks.
\textit{arXiv preprint arXiv:2005.11401}.

\bibitem{semanticscholar}
Kinney, R. M., Anastasiades, C., Authur, R., Beltagy, I., Bragg, J.,
Burachas, G., et al. (2023). The Semantic Scholar Open Data Platform.
\textit{arXiv preprint arXiv:2301.10140}.

\bibitem{asl2024fair}
Aghajani Asl, M., Asgari-Bidhendi, M., \& Minaei-Bidgoli, B. (2024).
FAIR-RAG: Faithful Adaptive Iterative Refinement for Retrieval-Augmented
Generation. \textit{arXiv preprint arXiv:2510.22344}.

\bibitem{keshav2007}
Keshav, S. (2007). How to Read a Paper. \textit{David R. Cheriton School of
Computer Science, University of Waterloo}.

\bibitem{godinez2025}
Godinez, A. (2025). HySemRAG: A Hybrid Semantic Retrieval-Augmented Generation
Framework for Automated Literature Synthesis and Methodological Gap Analysis.
\textit{arXiv preprint arXiv:2508.05666}.

\bibitem{mao2024}
Mao, Y., Dong, X., Xu, W., Gao, Y., Wei, B., \& Zhang, Y. (2024). FIT-RAG:
Black-Box RAG with Factual Information and Token Reduction. \textit{arXiv
preprint arXiv:2403.14374}.

\bibitem{badami2023}
Badami, M., Benatallah, B., \& Baez, M. (2023). Adaptive Search Query
Generation and Refinement in Systematic Literature Review. \textit{Information
Systems}, 102231.

\bibitem{yan2024}
Yan, S.-Q., Gu, J.-C., Zhu, Y., \& Ling, Z.-H. (2024). Corrective Retrieval
Augmented Generation. \textit{arXiv preprint arXiv:2401.15884}.

\bibitem{explainpaper}
Explainpaper. (2022). Explainpaper: Read papers faster.
\url{https://www.explainpaper.com}

\bibitem{qdrant2025}
Qdrant Team. (2025). Qdrant. \url{https://qdrant.tech/}

\bibitem{cormack2009}
Cormack, G. V., Clarke, C. L., \& Buettcher, S. (2009). Reciprocal rank fusion
outperforms Condorcet and individual rank learning methods. \textit{Proceedings
of the 32nd ACM SIGIR Conference}, 758--759.

\bibitem{es2023ragas}
Es, S., James, J., Espinosa-Anke, L., \& Schockaert, S. (2023). RAGAS:
Automated Evaluation of Retrieval Augmented Generation. \textit{arXiv preprint
arXiv:2309.15217}.

\bibitem{niutrans2021}
NiuTrans Research. (2021). ABigSurvey: A collection of survey papers on NLP and
AI subfields. \url{https://github.com/NiuTrans/ABigSurvey}

\bibitem{rahman2025automated}
Rahman, S., Shanto, H. K., Koana, U. A., \& Danish, S. M. (2025). Automated
Research Article Classification and Recommendation Using NLP and Machine
Learning. \textit{arXiv preprint arXiv:2510.05495}.


\bibitem{ghanadian2026}
Ghanadian, H., Kamali, A., \& Tekieh, M. H. (2026). Enhancing Scientific
Literature Chatbots with Retrieval-Augmented Generation: A Performance
Evaluation of Vector and Graph-Based Systems. \textit{arXiv preprint
arXiv:2602.17856}.

\bibitem{docling2024}
IBM Research. (2024). Docling: An Open-Source Document Conversion Tool.
\url{https://github.com/DS4SD/docling}

\bibitem{nassar2022}
Nassar, A., Livathinos, N., Lysak, M., \& Staar, P. (2022). TableFormer:
Table Structure Understanding with Transformers. \textit{Proceedings of the
IEEE/CVF Conference on Computer Vision and Pattern Recognition (CVPR)}.

\bibitem{shen2021}
Shen, Z., Zhang, R., Dell, M., Lee, B. C. G., Carlson, J., \& Li, W. (2021).
LayoutParser: A Unified Toolkit for Deep Learning Based Document Image
Analysis. \textit{Document Analysis and Recognition -- ICDAR 2021}.

\bibitem{smith2007}
Smith, R. (2007). An Overview of the Tesseract OCR Engine. \textit{Ninth
International Conference on Document Analysis and Recognition (ICDAR 2007)}.

\bibitem{du2020}
Du, Y., Chen, C., Jia, R., et al. (2020). PP-OCR: A Practical Ultra
Lightweight OCR System. \textit{arXiv preprint arXiv:2009.09941}.

\bibitem{councill2008}
Councill, I. G., Giles, C. L., \& Kan, M. Y. (2008). ParsCit: An Open-source
CRF Reference String Parsing Package. \textit{Proceedings of the Sixth
International Conference on Language Resources and Evaluation (LREC 2008)}.

\bibitem{lopez2009}
Lopez, P. (2009). GROBID: Combining Automatic Bibliographic Data Recognition
and Term Extraction for Scholarship Publications. \textit{Research and Advanced
Technology for Digital Libraries, ECDL 2009}.

\bibitem{tkaczyk2015}
Tkaczyk, D., Szostek, P., Fedoryszak, M., Dendek, P. J., \& Bolikowski, L.
(2015). CERMINE: Automatic Extraction of Structured Metadata from Scientific
Literature. \textit{International Journal on Document Analysis and Recognition,
18}(4), 317--335.

\bibitem{vila2021}
Vila, L., Subramanian, S., Lo, K., Beltagy, I., \& Downey, D. (2021). VILA:
Improving Structured Content Extraction from Scientific PDFs Using Visual
Layout Groups. \textit{Transactions of the Association for Computational
Linguistics, 9}, 1072--1088.

\bibitem{wei2022}
Wei, J., Tay, Y., Bommasani, R., et al. (2022). Emergent Abilities of Large
Language Models. \textit{Transactions on Machine Learning Research}.

\bibitem{li2020}
Li, M., Xu, Y., Cui, L., Huang, S., Wei, F., Li, Z., \& Zhou, M. (2020).
DocBank: A Benchmark Dataset for Document Layout Analysis. \textit{Proceedings
of the 28th International Conference on Computational Linguistics (COLING
2020)}.

\bibitem{keshav2007}
Keshav, S. (2007). How to Read a Paper. \textit{ACM SIGCOMM Computer
Communication Review, 37}(3), 83--84.

\bibitem{kintsch1978}
Kintsch, W., \& van Dijk, T. A. (1978). Toward a Model of Text Comprehension
and Production. \textit{Psychological Review, 85}(5), 363--394.

\bibitem{yang2023}
Lee, Y., Lee, K., Park, S., Hwang, D., Kim, J., \& Lee, H. (2023). QASA:
Advanced Question Answering on Scientific Articles. \textit{Proceedings of the
40th International Conference on Machine Learning (ICML 2023)}.

\bibitem{mao2003}
Mao, S., Rosenfeld, A., \& Kanungo, T. (2003). Document Structure Analysis
Algorithms: A Literature Survey. \textit{Proceedings of SPIE -- Document
Recognition and Retrieval X, 5010}.

\bibitem{besrour2025}
Besrour, I., He, J., Schreieder, T., \& Färber, M. (2025). SQuAI: Scientific
Question-Answering with Multi-Agent Retrieval-Augmented Generation.
\textit{arXiv preprint arXiv:2510.15682}.

\bibitem{es2023}
Es, S., James, J., Espinosa-Anke, L., \& Schockaert, S. (2023). Ragas:
Automated Evaluation of Retrieval Augmented Generation. \textit{arXiv preprint
arXiv:2309.15217}.

\bibitem{ragasdocs2025}
Ragas Documentation. (2025). Available Metrics.
\url{https://docs.ragas.io/en/stable/concepts/metrics/available_metrics/}

\bibitem{evidently2025}
Evidently AI. (2025). RAG Evaluation: Metrics, Testing and Best Practices.
\url{https://www.evidentlyai.com/llm-guide/rag-evaluation}

\bibitem{yu2024}
Yu, H., Gan, A., Zhang, K., Tong, S., Liu, Q., \& Liu, Z. (2024). Evaluation
of Retrieval-Augmented Generation: A Survey.
\url{https://doi.org/10.48550/arXiv.2405.07437}

\bibitem{shorten2026}
Shorten, C., Skaburskas, A., Jones, D. M., Pierse, C., Esposito, R.,
Trengrove, J., Dilocker, E., \& van Luijt, B. (2026). IRPAPERS: A Visual
Document Benchmark for Scientific Retrieval and Question Answering.
\textit{arXiv preprint arXiv:2602.17687}.

\bibitem{chen2026}
Chen, H., Long, Q., Pu, S., Luo, X., Ju, W., Xiao, M., Zhou, Y., Zhao, J.,
\& Wang, X. (2026). DeepEra: A Deep Evidence Reranking Agent for Scientific
Retrieval-Augmented Generated Question Answering. \textit{arXiv preprint
arXiv:2601.16478}.

\end{thebibliography}

% ============================================================
% APPENDICES
% ============================================================
\appendix

\chapter{API Endpoint Reference}
% TBD – Content to be added: list of FastAPI endpoints with request/response
% schemas and example calls.

\chapter{Database Schema}
% TBD – Content to be added: PostgreSQL table definitions, key foreign-key
% relationships, and an ER diagram if available.

\chapter{Example JSON Outputs}
% TBD – Content to be added: sample hierarchy JSON, chunk metadata payload, and
% guide generation output for a representative paper.

\chapter{Additional Figures and Screenshots}
% TBD – Content to be added: additional UI screenshots, LangGraph workflow
% diagrams, and any supplementary figures referenced in the main text.

\end{document}